% !TEX program = xelatex
% =====================================================================
%  Início do pré-natal e baixo peso ao nascer — SINASC, São Paulo, 2024
%  Relatório técnico-científico — Repasse acadêmico, Etapa 2
%  (a Etapa 1 permanece como registro histórico em PACOTE_REDACAO_ETAPA_1.pdf)
%
%  Compilador: XeLaTeX (Overleaf: Menu > Settings > Compiler > XeLaTeX).
%  Bibliografia: BibTeX + abntex2cite (ABNT, autor-data), arquivo referencias.bib.
%  Figuras: arquivos .tex em figuras/ (pgfplots/TikZ), editáveis como texto.
% =====================================================================
\documentclass[11pt,a4paper]{article}

% ---------- Identificação do documento (editar aqui) -----------------
\newcommand{\docTitulo}{Início do pré-natal e baixo peso ao nascer: análise observacional com dados do SINASC}
\newcommand{\docSubtitulo}{Material metodológico e resultados — São Paulo (capital), 2024, com série histórica 2014–2024}
\newcommand{\docEtapa}{Repasse acadêmico — Etapa 2}
\newcommand{\docCurso}{Mestrado em Computação Aplicada}
\newcommand{\docDisciplina}{Big Data \& Analytics}
\newcommand{\docLocal}{São Paulo}
\newcommand{\docAno}{2026}
% Autores e orientação: preencher entre as chaves (separar nomes com \\).
% Enquanto estiverem vazios, as linhas correspondentes não aparecem na capa.
\newcommand{\docAutores}{}
\newcommand{\docOrientacao}{}

% ---------- Motor e fontes --------------------------------------------
\usepackage{iftex}
\ifPDFTeX
  \errmessage{Este projeto usa fontes OpenType: selecione o compilador XeLaTeX
    (Overleaf: Menu > Settings > Compiler > XeLaTeX)}
\fi
\usepackage{amsmath}
\usepackage{fontspec}
\usepackage{unicode-math}
% Família Libertinus (incluída no TeX Live/Overleaf): serifada no texto,
% sem serifa em títulos e figuras, matemática compatível.
\setmainfont{LibertinusSerif}[Extension=.otf, UprightFont=*-Regular,
  BoldFont=*-Semibold, ItalicFont=*-Italic, BoldItalicFont=*-SemiboldItalic]
\setsansfont{LibertinusSans}[Extension=.otf, UprightFont=*-Regular,
  BoldFont=*-Bold, ItalicFont=*-Italic]
\setmonofont{LibertinusMono}[Extension=.otf, UprightFont=*-Regular, Scale=0.92]
\setmathfont{LibertinusMath-Regular.otf}
\usepackage[brazilian]{babel}

% ---------- Página ------------------------------------------------------
\usepackage[a4paper, top=2.5cm, bottom=2.5cm, left=2.5cm, right=2.5cm,
  headheight=22pt, headsep=0.8cm, footskip=1.1cm]{geometry}
\usepackage{setspace}
\setstretch{1.10}
\setlength{\parindent}{1.25em}
\setlength{\parskip}{0pt plus 1pt}
\usepackage{microtype}
\usepackage{etoolbox}
\usepackage{enumitem}
\setlist{itemsep=2pt, topsep=4pt, parsep=0pt}
\setlist[enumerate]{label=\arabic*., leftmargin=2.2em}
\setlist[itemize]{label=\textcolor{iptazul}{\small\textbullet}, leftmargin=1.6em}

% ---------- Cores ---------------------------------------------------------
% Cores amostradas do logotipo oficial do IPT (arquivo vetorial em ipt.br,
% página Comunicação Institucional). São valores medidos no arquivo,
% não códigos oficiais publicados em manual de marca.
\usepackage{xcolor}
\definecolor{iptazul}{HTML}{123C59}   % azul institucional escuro
\definecolor{iptciano}{HTML}{00B3E2}  % azul-ciano (pequenos destaques)
\definecolor{iptcinza}{HTML}{636362}  % cinza do logotipo
\colorlet{cinzaserie}{iptcinza!75}    % cinza médio para séries secundárias

% ---------- Tabelas -------------------------------------------------------
\usepackage{array}
\usepackage{booktabs}
\usepackage{colortbl}
\usepackage{tabularx}
\usepackage{siunitx}
\sisetup{output-decimal-marker={,}, group-separator={.}, group-minimum-digits=4,
  group-digits=integer, table-number-alignment=center}
\arrayrulecolor{iptazul}
\newcolumntype{L}{>{\raggedright\arraybackslash}X}
\newcommand{\cab}[1]{{\sffamily\color{iptazul}#1}}
\newcommand{\destaque}{\rowcolor{iptazul!6}}

% ---------- Figuras -----------------------------------------------------------
\usepackage{graphicx}
\usepackage{tikz}
\usetikzlibrary{arrows.meta, calc, positioning}
\usepackage{pgfplots}
\pgfplotsset{
  compat=1.18,
  estilo ipt/.style={
    axis lines*=left,
    axis line style={iptcinza!70, line width=0.5pt},
    tick style={iptcinza!70},
    tick align=outside,
    tick label style={font=\sffamily\footnotesize, text=iptcinza!30!black},
    label style={font=\sffamily\footnotesize, text=iptcinza!30!black},
    title style={font=\sffamily\footnotesize\bfseries, text=iptazul},
    major grid style={iptcinza!18, line width=0.4pt},
    legend style={font=\sffamily\footnotesize, draw=none, fill=white,
      fill opacity=0.9, text opacity=1},
    legend cell align=left,
    /pgf/number format/use comma,
    /pgf/number format/1000 sep={.},
  },
  serie principal/.style={iptazul, line width=1.2pt, mark size=2.2pt,
    mark options={fill=iptazul, solid}},
  serie secundaria/.style={cinzaserie, line width=1.2pt, mark size=2.2pt,
    mark options={fill=white, solid}},
}
\tikzset{rotulo valor/.style={font=\sffamily\scriptsize, text=iptcinza!20!black,
  inner sep=2.5pt}}

% ---------- Legendas (padrão ABNT: título acima, fonte abaixo) ---------------
\usepackage{caption}
\captionsetup{font=small, labelfont={sf, color=iptazul}, labelsep=endash,
  justification=justified, singlelinecheck=false, skip=5pt, position=top}
\newcommand{\fonte}[1]{\par\vspace{4pt}%
  {\footnotesize\color{iptcinza!50!black}\textsf{Fonte:} #1\par}}
\newcommand{\nota}[1]{\par\vspace{2pt}%
  {\footnotesize\color{iptcinza!50!black}\textsf{Nota:} #1\par}}

% ---------- Títulos ----------------------------------------------------------
\usepackage{titlesec}
\titleformat{\section}{\sffamily\Large\bfseries\color{iptazul}}{\thesection}{0.75em}{}
\titleformat{\subsection}{\sffamily\large\bfseries\color{iptazul}}{\thesubsection}{0.7em}{}
\titleformat{\subsubsection}{\sffamily\normalsize\bfseries\color{iptazul}}{\thesubsubsection}{0.6em}{}
\titleformat{\paragraph}[runin]{\sffamily\normalsize\bfseries\color{iptazul}}{}{0pt}{}
\titlespacing*{\paragraph}{0pt}{1.6ex plus 0.4ex minus 0.2ex}{0.6em}
\titlespacing*{\section}{0pt}{3.0ex plus 0.8ex minus 0.4ex}{1.4ex plus 0.3ex}
\titlespacing*{\subsection}{0pt}{2.4ex plus 0.6ex minus 0.3ex}{1.0ex plus 0.2ex}
\titlespacing*{\subsubsection}{0pt}{1.8ex plus 0.4ex minus 0.2ex}{0.6ex plus 0.2ex}
\setcounter{tocdepth}{1}

% Posicionamento de figuras e tabelas: evita páginas só com flutuantes
\renewcommand{\topfraction}{0.85}
\renewcommand{\bottomfraction}{0.6}
\renewcommand{\textfraction}{0.1}
\renewcommand{\floatpagefraction}{0.8}
\setlength{\textfloatsep}{14pt plus 4pt minus 4pt}
\setlength{\floatsep}{12pt plus 3pt minus 3pt}

% ---------- Cabeçalho e rodapé --------------------------------------------------
\usepackage{fancyhdr}
\pagestyle{fancy}
\fancyhf{}
\fancyhead[L]{\sffamily\footnotesize\color{iptcinza}IPT \textperiodcentered\ \docCurso}
\fancyhead[R]{\sffamily\footnotesize\color{iptcinza}\docEtapa}
\fancyfoot[R]{\sffamily\footnotesize\color{iptazul}\thepage}
\renewcommand{\headrulewidth}{0.5pt}
\renewcommand{\headrule}{\hbox to\headwidth{\color{iptciano}\leaders\hrule height \headrulewidth\hfill}}
\fancypagestyle{plain}{\fancyhf{}\fancyfoot[R]{\sffamily\footnotesize\color{iptazul}\thepage}%
  \renewcommand{\headrulewidth}{0pt}}

% ---------- Notação e atalhos --------------------------------------------------
\newcommand{\campo}[1]{\texttt{#1}}       % nomes de campos do SINASC
\newcommand{\pp}{\,p.p.}
\newcommand{\E}{\mathbb{E}}
\newcommand{\ATO}{\mathrm{ATO}}

% ---------- Hiperlinks e citações (abntex2cite depois do hyperref) --------------
\usepackage{hyperref}
\hypersetup{colorlinks=true, linkcolor=iptazul, citecolor=iptazul,
  urlcolor=iptazul, pdftitle={\docTitulo}, pdfsubject={\docSubtitulo},
  pdfkeywords={SINASC, pré-natal, baixo peso ao nascer, escore de propensão, ATO},
  pdflang={pt-BR}}
\urlstyle{same}
\usepackage[alf, abnt-etal-cite=3, abnt-etal-list=0]{abntex2cite}
% Caminhos de arquivos do projeto (sem os sinais < > que a ABNT usa em URLs)
\DeclareUrlCommand\arquivo{\urlstyle{tt}\def\UrlLeft{}\def\UrlRight{}}

% =====================================================================
\begin{document}

% ---------------------------------- Capa ---------------------------------------
\begin{titlepage}
\newgeometry{top=2.3cm, bottom=2.3cm, left=2.7cm, right=2.7cm}
\setlength{\parindent}{0pt}
\sffamily
% Logotipo oficial do IPT (PDF vetorial, sem alterações); o recorte apenas
% elimina a margem branca da página original.
\noindent\includegraphics[height=1.75cm, trim=192 203 192 203, clip]{logo_ipt.pdf}

\vspace{0.75cm}
{\small\color{iptcinza}\docCurso\\[1pt] Disciplina \docDisciplina\par}

\vspace*{\fill}
{\color{iptciano}\rule{2.4cm}{1.5pt}}\par
\vspace{0.7cm}
{\fontsize{23}{29}\selectfont\bfseries\color{iptazul}\raggedright\docTitulo\par}
\vspace{0.75cm}
{\rmfamily\Large\itshape\color{iptcinza!70!black}\raggedright\docSubtitulo\par}
\vspace{1.3cm}
{\small\color{iptazul}\addfontfeatures{LetterSpace=6}\MakeUppercase{\docEtapa}\par}
\ifdefempty{\docAutores}{}{%
  \vspace{1.1cm}{\normalsize\color{iptcinza!20!black}\docAutores\par}}
\ifdefempty{\docOrientacao}{}{%
  \vspace{0.5cm}{\small\color{iptcinza!20!black}\docOrientacao\par}}
\vspace*{\fill}

{\footnotesize\color{iptcinza}\begin{minipage}{0.78\linewidth}
Relatório técnico-científico elaborado no âmbito acadêmico do \docCurso{} do
Instituto de Pesquisas Tecnológicas. Não constitui publicação oficial da
instituição nem versão final de dissertação ou artigo.
\end{minipage}\par}
\vspace{0.6cm}
{\small\color{iptazul}\docLocal, \docAno\par}
\restoregeometry
\end{titlepage}

% ------------------------------- Sumário -----------------------------------------
\pagestyle{fancy}
\tableofcontents
\thispagestyle{fancy}

\section*{Apresentação}
\addcontentsline{toc}{section}{Apresentação}

Este relatório técnico-científico apresenta uma análise observacional da relação entre o
momento de início do pré-natal e o baixo peso ao nascer no município de São Paulo. A
pergunta principal é se, entre nascidos vivos de mães residentes na capital em 2024, o
início do pré-natal nos meses 1--3 está associado a menor risco de baixo peso ao nascer
(peso inferior a 2.500 g) do que o início nos meses 4--6, após ajuste por características
maternas observadas.

Os dados são os microdados públicos do Sistema de Informações sobre Nascidos Vivos
(SINASC), com um registro por Declaração de Nascido Vivo. O município de São Paulo foi
escolhido pela qualidade de preenchimento das variáveis essenciais, critério definido sem
uso da associação com o desfecho (Seção~\ref{sec:qualidade}). A população analítica reúne
114.694 nascimentos de 2024, de gestação única, com peso entre 500 e 6.000 g e início do
pré-natal nos meses 1--6: 104.576 com início nos meses 1--3 ($T=1$) e 10.118 nos meses
4--6 ($T=0$) (Seção~\ref{sec:populacao}).

As diferenças entre os grupos nas características maternas observadas foram tratadas com
escore de propensão logístico e pesos de sobreposição (ATO), que estimam a diferença de
risco na população em que os dois momentos de início são frequentes
(Seção~\ref{sec:ato}). A incerteza foi quantificada por bootstrap não paramétrico com
1.000 réplicas, com reestimação do escore em cada réplica (Seção~\ref{sec:incerteza}).

Em 2024, o risco ponderado de baixo peso foi de 8,68\% em $T=1$ e de 10,42\% em $T=0$,
diferença de $-1{,}74$\pp{} (IC 95\%: $-2{,}36$ a $-1{,}12$\pp;
Seção~\ref{sec:resultados}). A mesma especificação, repetida separadamente em cada ano de
2014 a 2024, produziu diferenças negativas nos onze anos, apresentadas como estimativas
pontuais, sem intervalos de confiança anuais (Seção~\ref{sec:serie}). Uma análise
descritiva da duração gestacional e da seleção da população examina a temporalidade do
início do pré-natal; ela caracteriza a composição dos grupos, mas não identifica a
magnitude nem a direção de um possível viés (Seção~\ref{sec:selecao}).

O resultado principal é uma associação ajustada por características maternas observadas, e
não um efeito causal identificado. Sua leitura causal dependeria de hipóteses não
verificáveis com o SINASC, entre elas a ausência de confundimento não medido; permanecem
também em aberto a temporalidade da exposição, a seleção por nascidos vivos, a dependência
entre observações e o alcance restrito da população de sobreposição
(Seção~\ref{sec:discussao}). Essa distinção orienta todo o texto.

Esta é a segunda etapa do relatório. Em relação à primeira, restrita à estimação de 2024,
acrescenta a série 2014--2024 e a análise de temporalidade e seleção; a pergunta, a
população, o estimando, o método e os resultados de 2024 não mudaram. Uma exploração
descritiva preliminar com dados do Sistema de Informações sobre Mortalidade (SIM),
discutida na versão anterior, foi retirada do escopo: o estudo usa exclusivamente o SINASC
e não realizou vinculação individual entre os dois sistemas.

Nenhuma análise foi reexecutada para este relatório. Os números foram transcritos das
saídas salvas dos notebooks do projeto e conferidos com elas, e as figuras foram
redesenhadas a partir dos mesmos valores que geraram os gráficos dos notebooks.

Convenções: diferenças calculadas sempre como $T=1$ menos $T=0$; p.p.\ = pontos
percentuais; SMD = diferença média padronizada; IC 95\% = intervalo de confiança de 95\%;
ATO = efeito médio na população de sobreposição (\emph{average treatment effect in the
overlap population}). Abaixo de cada tabela e figura indica-se o notebook, e a seção, de
onde os valores foram extraídos.

% =====================================================================
\section{Introdução}\label{sec:introducao}

\subsection{Baixo peso ao nascer}

A Organização Mundial da Saúde define baixo peso ao nascer como peso inferior a
2.500\,g \cite{who_lbw}. A definição não depende da idade gestacional: o baixo peso pode
resultar de restrição do crescimento intrauterino, de prematuridade ou de ambas
\cite{who_lbw}. O indicador agrega, portanto, mecanismos distintos, o que tem
consequências diretas para a interpretação de qualquer associação com fatores da
assistência pré-natal (Seção~\ref{sec:discussao}).

A relevância do desfecho para a saúde materno-infantil é bem estabelecida. Segundo a OMS,
recém-nascidos com baixo peso têm probabilidade de morrer cerca de vinte vezes maior que
recém-nascidos mais pesados, e o indicador integra a lista de referência global de
indicadores básicos de saúde da organização \cite{who_lbw}. No Brasil, o peso ao nascer é
registrado para todos os nascidos vivos na Declaração de Nascido Vivo, o que permite
estudá-lo em base populacional.

\subsection{O momento de início do pré-natal}

O início precoce do pré-natal permite rastreamento, suplementação e manejo antecipado de
condições de risco. Essa é a hipótese substantiva que motivou o projeto: quando o
acompanhamento começa no primeiro trimestre, haveria mais oportunidade de intervir sobre
fatores associados à prematuridade e à restrição de crescimento. A assistência pré-natal
no Brasil foi avaliada em âmbito nacional por \citeonline{viellas2014}, e estudos com
dados do SINASC associaram índices de adequação do pré-natal ao baixo peso
\cite{vale2021}. Esses índices, contudo, combinam o mês de início com o número de
consultas; o contraste estudado aqui isola o momento de início.

A evidência sobre efeitos do pré-natal em desfechos do nascimento é heterogênea, e
\citeonline{corman2018} discutem as dificuldades de identificá-los com dados
observacionais. Trabalhos recentes procuram enfrentar esse problema com métodos de
inferência causal, como a abordagem de decomposição causal aplicada por
\citeonline{vachuska2025} ao início precoce do pré-natal e às desigualdades raciais em
desfechos do nascimento.

\subsection{Diferenças entre gestantes e confundimento}

O momento de início do pré-natal não é atribuído ao acaso. Nos dados analisados, o início
nos meses 4--6 é mais frequente entre mães mais jovens, com menor escolaridade e sem união
formal registrada (Seção~\ref{sec:comparabilidade}). Características socioeconômicas,
reprodutivas e clínicas podem influenciar simultaneamente o momento de início e o peso ao
nascer; em estudo brasileiro de base populacional, escolaridade, raça/cor, idade e
paridade associaram-se ao baixo peso a termo \cite{falcao2020}. A comparação bruta entre
os grupos mistura, portanto, a eventual contribuição do momento de início com diferenças
de composição. Parte dessas diferenças é registrada no SINASC e pode ser ajustada; outra
parte --- renda, cobertura por plano de saúde, planejamento da gestação, tabagismo,
condições clínicas prévias --- não é registrada.

\subsection{Objetivo e pergunta de pesquisa}

O objetivo desta etapa é estimar a diferença no risco de baixo peso ao nascer entre
nascimentos cujas mães iniciaram o pré-natal nos meses 1--3 e nascimentos com início nos
meses 4--6, após ponderação por características maternas observadas, e documentar as
condições sob as quais essa estimativa pode ser interpretada.

\begin{center}
\begin{minipage}{0.9\linewidth}
\setstretch{1.08}
{\color{iptciano}\rule{\linewidth}{0.6pt}}\par\vspace{2pt}
\textsf{\color{iptazul}Pergunta principal.} Entre nascidos vivos de mães residentes no
município de São Paulo em 2024, o início do pré-natal nos meses 1--3 está associado a
menor risco de baixo peso ao nascer, em comparação com o início nos meses 4--6, após
ajuste por características maternas observadas?\par\vspace{2pt}
{\color{iptciano}\rule{\linewidth}{0.6pt}}
\end{minipage}
\end{center}

Convém distinguir duas perguntas. A pergunta causal que motivou o projeto é
contrafactual: qual seria a diferença no risco de baixo peso se as mesmas gestações
tivessem iniciado o pré-natal nos meses 1--3, e não nos meses 4--6? Respondê-la exige
hipóteses de identificação --- ausência de confundimento não medido, positividade e
consistência da exposição \cite{hernan2025} --- que os dados do SINASC não permitem
verificar. A pergunta efetivamente respondida nesta etapa é mais restrita: qual é a
diferença de risco entre os grupos depois que suas distribuições de características
maternas observadas são equilibradas por ponderação? O resultado é uma associação
ajustada, válida para uma população-alvo definida pelos próprios pesos
(Seção~\ref{sec:ato}).

Duas análises complementares, ambas descritivas, acompanham a estimação principal: a
repetição da mesma especificação em cada ano de 2014 a 2024 (Seção~\ref{sec:serie}) e o
exame da relação entre duração da gestação, oportunidade de início do pré-natal e seleção
da população analítica (Seção~\ref{sec:selecao}). Nenhuma delas altera a pergunta, a
população, o estimando ou o resultado da análise principal.

% =====================================================================
\section{Dados e população de origem}\label{sec:dados}

\subsection{O Sistema de Informações sobre Nascidos Vivos}

O Sistema de Informações sobre Nascidos Vivos (SINASC), do Ministério da Saúde, reúne as
informações da Declaração de Nascido Vivo, preenchida para cada nascimento vivo. Os
microdados públicos, disseminados pelo DATASUS, contêm características da mãe, da
gestação, do parto e do recém-nascido. A cobertura e a qualidade do sistema foram objeto
de avaliação nacional \cite{szwarcwald2019}. Em maternidades públicas da capital paulista,
\citeonline{bonilha2018} encontraram boa concordância entre o SINASC e a coleta de campo
para peso ao nascer e tipo de gravidez, e concordância menor para campos de assistência e
de história reprodutiva.

A unidade de análise é o registro de nascido vivo. Cada linha corresponde a uma
declaração, e não a uma gestante: a mesma mulher pode contribuir com mais de um registro
em gestações múltiplas (excluídas da população analítica) ou em mais de um parto no mesmo
ano. Os dados não formam uma coorte de gestantes acompanhadas desde o início da gestação;
exposição, covariáveis e desfecho são registrados no momento do parto.

O projeto dispõe de bases locais completas do SINASC para os onze anos de 2014 a 2024. As
definições de campos seguem o dicionário analítico do projeto, que corresponde ao layout
documentado até 2019; códigos não enumerados nesse dicionário, como limites de idade
materna e sentinelas do campo de perdas fetais, não foram presumidos.

\subsection{Qualidade dos dados e escolha da capital}\label{sec:qualidade}

A orientação metodológica foi concentrar a investigação em localidades com melhor
qualidade de registro. O recorte foi escolhido exclusivamente por critérios de
preenchimento, validade e retenção de registros, \emph{sem consulta à magnitude da
associação entre início do pré-natal e baixo peso}. Dois levantamentos fundamentaram a
decisão.

O primeiro comparou seis recortes geográficos em 2024 (Tabela~\ref{tab:recortes}). A
capital apresentou a maior proporção de registros com as três variáveis essenciais
simultaneamente válidas --- mês de início do pré-natal (\campo{MESPRENAT} entre 1 e 9),
tipo de gravidez (\campo{GRAVIDEZ} entre 1 e 3) e peso inteiro entre 500 e 6.000\,g:
99,04\%, contra 97,93\% no estado de São Paulo e 97,16\% no Sudeste. O mês de início do
pré-natal, que define a exposição, foi válido em 99,15\% dos registros da capital,
98,08\% do estado e 97,32\% do Sudeste. Escolaridade, raça/cor, situação conjugal e
perdas fetais anteriores também tiveram preenchimento maior na capital que no estado.

\begin{table}[htbp]
\centering
\caption{Qualidade e volume de registros por recorte geográfico de residência materna, SINASC 2024}
\label{tab:recortes}
\small
\begin{tabular}{@{}l S[table-format=6.0] S[table-format=2.2] S[table-format=6.0]
  S[table-format=2.2] S[table-format=6.0] S[table-format=5.0]@{}}
\toprule
\cab{Recorte} & {\cab{Nascimentos}} & {\cab{Essenciais (\%)}} & {\cab{Elegíveis}} &
  {\cab{Retenção (\%)}} & {\cab{$T=1$}} & {\cab{$T=0$}} \\
\midrule
Sudeste                 & 904365 & 97.16 & 842998 & 93.21 & 758924 & 84074 \\
São Paulo --- estado    & 470887 & 97.93 & 442733 & 94.02 & 402207 & 40526 \\
\destaque
São Paulo --- capital   & 120789 & 99.04 & 114694 & 94.95 & 104576 & 10118 \\
Minas Gerais            & 219819 & 95.49 & 201200 & 91.53 & 182626 & 18574 \\
Rio de Janeiro          & 163722 & 96.93 & 151969 & 92.82 & 131899 & 20070 \\
Espírito Santo          &  49937 & 98.06 &  47096 & 94.31 &  42192 &  4904 \\
\bottomrule
\end{tabular}
\nota{Essenciais: \campo{MESPRENAT} 1--9, \campo{GRAVIDEZ} 1--3 e peso inteiro de 500 a
6.000\,g simultaneamente. Elegíveis: gestação única, peso de 500 a 6.000\,g e início do
pré-natal nos meses 1--6, aplicados a cada recorte. Retenção: elegíveis sobre nascimentos. Os recortes se
sobrepõem (capital $\subset$ estado $\subset$ Sudeste) e não devem ser somados.}
\fonte{\arquivo{outputs/qualidade_geografica_2024/RECOMENDACAO_RECORTE.md} e
\arquivo{comparacao_recortes_2024.csv}, na mesma pasta.}
\end{table}

O segundo levantamento examinou a série 2014--2024 para verificar se a vantagem da capital
era persistente (Figura~\ref{fig:qualidade}). A validade de \campo{MESPRENAT} na capital
variou de 98,45\% em 2014 a 99,15\% em 2024 e foi superior à do restante do estado em
todos os onze anos, com vantagem entre 0,81 e 4,62\pp{} No período, a capital somou
1.587.749 nascimentos elegíveis segundo os critérios da Seção~\ref{sec:populacao} ---
1.401.494 com início nos meses 1--3 e 186.255 nos meses 4--6 ---, com retenção acumulada
de 94,61\% dos registros.

\begin{figure}[htbp]
\centering
\caption{Proporção de registros com mês de início do pré-natal válido, por ano de
nascimento e residência materna, estado de São Paulo, 2014--2024}
\label{fig:qualidade}
% Figura 1 — Proporção de registros com MESPRENAT válido (códigos 1–9), por ano.
% Valores copiados do gráfico Plotly salvo em
% notebooks/principal/05a_qualidade_geografica_sp_2014_2024.ipynb (seção 7).
% Denominador: todos os nascimentos do ano no recorte de residência materna.
\begin{tikzpicture}
\begin{axis}[
  estilo ipt,
  width=0.94\linewidth, height=6.1cm,
  xmin=2013.6, xmax=2024.4,
  ymin=93, ymax=100,
  xtick={2014,...,2024},
  xticklabel style={/pgf/number format/1000 sep={}},
  ytick={93,94,...,100},
  ymajorgrids,
  xlabel={Ano de nascimento},
  ylabel={MESPRENAT válido (\%)},
  legend style={at={(0.98,0.05)}, anchor=south east},
]
\addplot[serie principal, mark=*] coordinates {
  (2014,98.4480) (2015,98.4839) (2016,98.7980) (2017,99.0089) (2018,99.0324)
  (2019,98.9520) (2020,98.8038) (2021,98.9044) (2022,99.0966) (2023,99.1447)
  (2024,99.1547)};
\addlegendentry{Município de São Paulo (capital)}
\addplot[serie secundaria, mark=square*] coordinates {
  (2014,93.8857) (2015,93.8649) (2016,95.9476) (2017,97.2502) (2018,97.7259)
  (2019,97.9148) (2020,97.6013) (2021,97.9123) (2022,98.0692) (2023,98.3390)
  (2024,97.7061)};
\addlegendentry{Restante do estado de São Paulo}
% Rótulos seletivos: primeiro e último ano de cada série
\node[rotulo valor, anchor=south] at (axis cs:2014,98.4480) {98,45};
\node[rotulo valor, anchor=south] at (axis cs:2024,99.1547) {99,15};
\node[rotulo valor, anchor=north, yshift=-3pt] at (axis cs:2014,93.8857) {93,89};
\node[rotulo valor, anchor=north, yshift=-3pt] at (axis cs:2024,97.7061) {97,71};\end{axis}
\end{tikzpicture}

\nota{Denominador: todos os nascimentos do ano no recorte. Restante do estado: residência
no estado de São Paulo fora da capital.}
\fonte{SINASC 2014--2024;
\arquivo{notebooks/principal/05a_qualidade_geografica_sp_2014_2024.ipynb} (seção 7;
indicadores anuais da tabela \texttt{dados}).}
\end{figure}

A série também revelou um ponto de atenção. Na capital, o número de nascimentos com
início nos meses 4--6 caiu de 24.348 em 2014 (14,76\% dos elegíveis) para 10.118 em 2024
(8,82\%), redução de 58,44\%, enquanto o total de elegíveis caiu 30,48\%. A redução do
grupo de comparação não constitui, por si só, violação de positividade, mas motivou o
diagnóstico de comparabilidade da Seção~\ref{sec:comparabilidade}. As causas dessa
mudança não foram investigadas.

A estimação principal, com quantificação de incerteza, foi restrita a 2024. Os onze anos
serviram, primeiro, para avaliar qualidade e viabilidade amostral; depois, a mesma
especificação foi estimada separadamente em cada ano, de forma descritiva e sem intervalos
de confiança (Seção~\ref{sec:serie}). O estado de São Paulo foi registrado como
alternativa de maior abrangência, mas não foi analisado.

% =====================================================================
\section{População analítica}\label{sec:populacao}

\subsection{Critérios de seleção}

A população analítica foi definida sobre o SINASC 2024 com a projeção de seleção
congelada no código do projeto, acrescida dos filtros de município e de mês de início:

\begin{enumerate}
\item residência materna no município de São Paulo (\campo{CODMUNRES}~=~355030);
\item nascimento ocorrido em 2024;
\item gestação única (\campo{GRAVIDEZ}~=~1);
\item peso ao nascer inteiro entre 500 e 6.000\,g, faixa operacional do projeto e não
domínio clínico validado;
\item início do pré-natal conhecido e situado nos meses 1 a 6 (\campo{MESPRENAT} entre 1
e 6).
\end{enumerate}

Registros com covariáveis incompletas não foram excluídos. Foram verificadas a unicidade
e a ausência de nulos na chave do registro (\campo{CONTADOR}), a ocorrência de todas as
datas de nascimento em 2024 e a igualdade das contagens com a referência do levantamento
de qualidade.

\subsection{Exposição e desfecho}

A exposição e o desfecho foram definidos para cada registro $i$ como
\begin{equation}\label{eq:def}
T_i =
\begin{cases}
1, & \text{se } \campo{MESPRENAT}_i \in \{1,2,3\},\\[2pt]
0, & \text{se } \campo{MESPRENAT}_i \in \{4,5,6\},
\end{cases}
\qquad
Y_i =
\begin{cases}
1, & \text{se o peso ao nascer é inferior a 2.500\,g},\\[2pt]
0, & \text{caso contrário.}
\end{cases}
\end{equation}
Por convenção, $T=1$ é chamado grupo de início precoce e $T=0$, grupo de comparação. Na
população analítica, os inícios distribuíram-se em 35.224 no mês 1, 54.444 no mês 2,
14.908 no mês 3, 5.366 no mês 4, 3.005 no mês 5 e 1.747 no mês 6. Inícios nos meses 7--9
ficaram fora do contraste: a ocorrência de inícios muito tardios depende ainda mais da
duração da gestação, problema discutido nas Seções~\ref{sec:selecao}
e~\ref{sec:temporalidade}.

\subsection{Covariáveis}\label{sec:covariaveis}

Todas as covariáveis são registradas no parto e foram tratadas como aproximações de
condições anteriores à exposição (Tabela~\ref{tab:variaveis}). O valor ignorado (código 9
ou nulo) foi mantido como indicador de ausência, nunca interpretado como categoria
substantiva.

\begin{table}[htbp]
\centering
\caption{Variáveis da análise, codificação e papel no modelo principal}
\label{tab:variaveis}
\small
\begin{tabularx}{\linewidth}{@{}>{\raggedright\arraybackslash}p{2.0cm}
  >{\raggedright\arraybackslash}p{3.3cm} >{\raggedright\arraybackslash}p{2.6cm} L@{}}
\toprule
\cab{Papel} & \cab{Variável (campo)} & \cab{Codificação} & \cab{Uso} \\
\midrule
Exposição & Início do pré-natal (\campo{MESPRENAT}) & 1--3: $T=1$; 4--6: $T=0$ &
  Define os grupos \\
Desfecho & Peso ao nascer (\campo{PESO}) & $<$ 2.500\,g: $Y=1$ &
  Não usado no ajuste nem no diagnóstico \\
\addlinespace
Covariável & Idade materna (\campo{IDADEMAE}) & Anos (10--59) e idade$^2$ &
  Incluída; o termo quadrático acomoda o desequilíbrio concentrado nas adolescentes \\
Covariável & Escolaridade (\campo{ESCMAE2010}) & 0--5; ignorado & Incluída \\
Covariável & Raça/cor (\campo{RACACORMAE}) & 1--5; ignorado & Incluída \\
Covariável & Paridade (\campo{PARIDADE}) & Nulípara; multípara & Incluída \\
Covariável & Perdas fetais anteriores (\campo{QTDFILMORT}) & 0; 1; 2 ou mais; ignorado &
  Incluída, com ressalvas de mensuração \\
\addlinespace
Diagnóstico & Situação conjugal (\campo{ESTCIVMAE}) & 1--5; ignorado &
  Excluída do modelo principal; usada no diagnóstico e na sensibilidade \\
\addlinespace
Vedadas & Idade gestacional, consultas, tipo de parto, Apgar, índice de Kotelchuck, grupo
  de Robson & --- & Posteriores à exposição ou mediadoras; nunca usadas como controle \\
\bottomrule
\end{tabularx}
\fonte{\arquivo{notebooks/principal/05b_comparabilidade_sp_capital_2024.ipynb} (seção 3);
\arquivo{notebooks/principal/06_ato_sp_capital_2024.ipynb} (seção 1).}
\end{table}

Algumas decisões merecem registro. O campo de perdas fetais anteriores é definido no
dicionário como número de perdas fetais e abortos anteriores à gestação atual, portanto
anterior à exposição; na população analisada, 99,98\% dos registros têm valor não
sentinela e 99,99\% estão entre 0 e 9, sem ocorrência do código 99. O campo mistura,
porém, perdas espontâneas e abortos, estes provavelmente sub-registrados, e a categoria
``2 ou mais'' limita a influência de valores altos pouco plausíveis. A situação conjugal
foi excluída do modelo principal por temporalidade frágil: casamento ou união podem
ocorrer depois do início do pré-natal. A idade não apresentou registros fora de 10--59
anos e não houve imputação. Outros campos possivelmente anteriores à exposição
(\campo{QTDGESTANT}, \campo{QTDFILVIVO}, \campo{IDADEPAI}, \campo{CODOCUPMAE}) não foram
incorporados nesta versão. O bloqueio das variáveis vedadas é verificado por código antes
do ajuste do escore; ajustar por mediadores ou por consequências da exposição introduziria
viés de sobreajuste \cite{schisterman2009}.

\subsection{Registros fora da população}

Dos 120.789 nascimentos de mães residentes na capital em 2024, 114.694 (94,95\%)
compuseram a população analítica. As exclusões foram aplicadas em sequência
(Tabela~\ref{tab:fluxo}) e têm naturezas distintas. Apenas a primeira etapa --- registros
sem informação válida de mês de início, tipo de gravidez ou peso --- está relacionada à
qualidade do registro. As demais são critérios de definição da população: gestações
múltiplas e inícios nos meses 7--9 são categorias válidas que ficam fora do contraste, e
o peso fora da faixa de 500 a 6.000\,g é uma exclusão operacional, não um erro presumido.
Nenhum registro foi excluído por covariável incompleta ou por falta de suporte no escore
de propensão.

\begin{table}[htbp]
\centering
\caption{Composição da população analítica a partir dos nascimentos de mães residentes no
município de São Paulo, 2024}
\label{tab:fluxo}
\small
\begin{tabularx}{\linewidth}{@{}L S[table-format=-6.0] S[table-format=3.2]
  >{\raggedright\arraybackslash}p{4.4cm}@{}}
\toprule
\cab{Etapa} & {\cab{Registros}} & {\cab{\% do total}} & \cab{Natureza} \\
\midrule
Nascimentos de mães residentes na capital & 120789 & 100.00 & --- \\
\addlinespace
\quad Sem \campo{MESPRENAT}, \campo{GRAVIDEZ} ou peso válidos & -1028 & 0.85 &
  Qualidade do registro \\
\quad Gestação múltipla & -3210 & 2.66 & Categoria válida fora da população \\
\quad Gestação única com peso fora de 500--6.000\,g & -122 & 0.10 &
  Exclusão operacional \\
\quad Gestação única com início nos meses 7--9 & -1735 & 1.44 &
  Categoria válida fora do contraste \\
\addlinespace
\destaque
População analítica & 114694 & 94.95 & $T=1$: 104.576; $T=0$: 10.118 \\
\bottomrule
\end{tabularx}
\nota{Etapas aplicadas em sequência; cada linha conta apenas os registros não excluídos
nas etapas anteriores. Percentuais sobre os 120.789 nascimentos.}
\fonte{\arquivo{outputs/qualidade_geografica_2024/comparacao_recortes_2024.csv} (recorte
capital, 2024); população analítica e inícios nos meses 7--9 conferidos em
\arquivo{notebooks/principal/08_temporalidade_prenatal_sp_capital_2024.ipynb} (seção 5.1).}
\end{table}

% =====================================================================
\section{Comparabilidade dos grupos e confundimento}\label{sec:comparabilidade}

\subsection{Por que a comparação bruta pode ser confundida}

Seja $X$ o vetor de características maternas registradas. Se $X$ influencia tanto a
probabilidade de início precoce quanto o risco de baixo peso, a diferença bruta de riscos
entre $T=1$ e $T=0$ combina dois componentes que não se separam sem hipóteses adicionais:
a eventual relação do momento de início com o desfecho e as diferenças de composição dos
grupos. A estrutura assumida está representada na Figura~\ref{fig:dag}
(Seção~\ref{sec:discussao}): as características observadas foram ajustadas; fatores não
registrados permanecem como fonte possível de confundimento; e variáveis posteriores ao
início do pré-natal --- número de consultas, idade gestacional, tipo de parto --- não
foram usadas como controle, pois condicionar nelas bloquearia parte da própria relação
estudada ou introduziria viés \cite{schisterman2009,wilcox2011}. O diagnóstico de
comparabilidade foi realizado antes de qualquer estimação e sem uso do desfecho.

\subsection{Diferença média padronizada}

O desequilíbrio entre os grupos foi descrito pela diferença média padronizada (SMD), que
expressa a diferença em unidades de desvio-padrão e é recomendada para diagnosticar
balanço em análises com escore de propensão \cite{austin2015}. Para uma covariável
numérica,
\begin{equation}\label{eq:smd-num}
\mathrm{SMD} = \frac{\bar{x}_1 - \bar{x}_0}{\sqrt{\left(s_1^2 + s_0^2\right)/2}},
\end{equation}
em que $\bar{x}_t$ e $s_t^2$ são a média e a variância amostral no grupo $t$. Para
covariáveis categóricas, cada nível é tratado como indicador binário com proporção $p_t$
no grupo $t$:
\begin{equation}\label{eq:smd-cat}
\mathrm{SMD} = \frac{p_1 - p_0}{\sqrt{\left[p_1(1-p_1) + p_0(1-p_0)\right]/2}},
\end{equation}
com $p_t$ calculado sobre todos os registros do grupo, inclusive a categoria ignorado. O
resumo de cada variável é o maior $|\mathrm{SMD}|$ entre seus níveis. Os valores 0,1 e
0,2 foram usados como referências convencionais, não como testes. P-valores não foram
usados como critério de equilíbrio: com mais de cem mil registros, diferenças triviais
seriam estatisticamente significativas.

\begin{table}[htbp]
\centering
\caption{Desequilíbrio das covariáveis antes de qualquer ajuste: maior $|\mathrm{SMD}|$
por variável, São Paulo (capital), 2024}
\label{tab:smd}
\small
\begin{tabular}{@{}l S[table-format=1.3] l S[table-format=2.1] S[table-format=2.1]@{}}
\toprule
\cab{Variável} & {\cab{$|\mathrm{SMD}|$}} & \cab{Nível mais desequilibrado} &
  {\cab{$T=1$}} & {\cab{$T=0$}} \\
\midrule
Escolaridade             & 0.452 & Superior completo (\%) & 31.1 & 12.9 \\
Situação conjugal\textsuperscript{a} & 0.421 & Casada (\%) & 36.5 & 18.1 \\
Idade materna            & 0.270 & Média (anos)           & 29.3 & 27.5 \\
Raça/cor                 & 0.199 & Branca (\%)            & 45.9 & 36.1 \\
Paridade                 & 0.166 & Nulípara (\%)          & 40.4 & 32.4 \\
Perdas fetais anteriores & 0.041 & Uma (\%)               & 17.1 & 15.6 \\
\bottomrule
\end{tabular}
\nota{\textsuperscript{a} Não incluída no modelo principal (temporalidade frágil). SMD
calculado sem uso do desfecho.}
\fonte{\arquivo{notebooks/principal/05b_comparabilidade_sp_capital_2024.ipynb} (seção 4;
tabela \texttt{saida\_balanco}).}
\end{table}

Escolaridade e situação conjugal concentraram os maiores desequilíbrios
(Tabela~\ref{tab:smd}). A idade média foi 29,3 anos em $T=1$ e 27,5 em $T=0$, e mães de
10--19 anos foram mais frequentes em $T=0$ (11,9\% contra 6,6\%). Raça/cor e paridade
apresentaram desequilíbrio moderado; perdas fetais anteriores, desequilíbrio pequeno. O
padrão é compatível com um gradiente socioeconômico: o início nos meses 4--6 é mais
frequente entre mães mais jovens, menos escolarizadas e sem união formal registrada.

\subsection{Escore de propensão e sobreposição}\label{sec:overlap}

O escore de propensão é a probabilidade condicional de pertencer ao grupo $T=1$ dadas as
covariáveis,
\begin{equation}\label{eq:ps}
e(X) = P(T = 1 \mid X),
\end{equation}
e resume em uma dimensão as diferenças de composição entre os grupos nas covariáveis
incluídas. A hipótese de positividade requer $0 < e(x) < 1$ para os perfis da
população-alvo: cada combinação de características deve ter probabilidade não nula de
pertencer a ambos os grupos \cite{petersen2012}. Empiricamente, examina-se a sobreposição
(suporte comum) das distribuições do escore estimado em $T=1$ e $T=0$. Quando a
sobreposição é limitada, estimadores ponderados podem depender de poucos registros com
pesos elevados, e restringir a análise à região de sobreposição altera a população para a
qual o resultado vale \cite{crump2009,petersen2012}.

No diagnóstico, o escore foi estimado por regressão logística com penalização L2 e seis
covariáveis --- incluindo situação conjugal e sem o termo quadrático de idade ---, com
predições fora da partição de treino (validação cruzada estratificada em cinco partições).
Os principais resultados foram:
\begin{itemize}
\item AUC fora da partição de 0,655, discriminação modesta: as covariáveis registradas
explicam pouco do momento de início do pré-natal;
\item escores entre 0,514 e 0,986 em $T=1$ e entre 0,712 e 0,981 em $T=0$, sem registros
abaixo de 0,05 ou acima de 0,99;
\item 125 registros de $T=1$ e nenhum de $T=0$ fora do suporte comum observado;
\item cauda superior com poucos controles: acima de 0,95 estavam 28.683 registros de $T=1$
(27,4\% do grupo) e 982 de $T=0$ (9,7\%); entre 0,975 e 1,000, 5.366 contra 106.
\end{itemize}

A AUC modesta favorece a sobreposição, mas não indica validade causal. Ela mostra apenas
que as covariáveis observadas discriminam pouco os grupos, o que é compatível com a
existência de determinantes importantes do momento de início que não estão registrados. O
escore descreve as covariáveis incluídas e nada informa sobre as não observadas.

O diagnóstico revelou uma combinação específica: suporte comum amplo, mas escores
concentrados perto de 1, reflexo de $P(T=1) \approx 0{,}91$. Nessas condições, o recorte
simétrico convencional de 0,05 a 0,95 retiraria 27,4\% dos registros de $T=1$, e pesos
de probabilidade inversa seriam instáveis: sob ponderação voltada aos tratados, controles
com escore alto receberiam pesos proporcionais à razão $e/(1-e)$, que chega a cerca de 52.
Essa combinação motivou a escolha da ponderação por sobreposição.

% =====================================================================
\section{Metodologia: ponderação por sobreposição (ATO)}\label{sec:ato}

\subsection{Estimandos: ATE, ATT e ATO}

Na notação de resultados potenciais, $Y_i(1)$ e $Y_i(0)$ denotam o desfecho que o
nascimento $i$ apresentaria sob início do pré-natal nos meses 1--3 e nos meses 4--6,
respectivamente; apenas um deles é observado. Seja
$\tau(x) = \E[\,Y(1) - Y(0) \mid X = x\,]$ a diferença média condicional ao perfil $x$. Os
estimandos usuais diferem na população sobre a qual $\tau(x)$ é promediado:
\begin{align}
\tau_{\mathrm{ATE}} &= \E\left[\,Y(1) - Y(0)\,\right], \label{eq:ate}\\
\tau_{\mathrm{ATT}} &= \E\left[\,Y(1) - Y(0) \mid T = 1\,\right], \label{eq:att}\\
\tau_{\ATO} &= \frac{\E\left[\,h(X)\,\tau(X)\,\right]}{\E\left[\,h(X)\,\right]},
  \qquad h(x) = e(x)\,\left[1 - e(x)\right]. \label{eq:ato}
\end{align}
O ATE refere-se a toda a população; o ATT, aos nascimentos do grupo $T=1$. O ATO
refere-se a uma população-alvo cuja densidade é proporcional a $f(x)\,e(x)\,[1-e(x)]$, em
que $f$ é a densidade das covariáveis \cite{li2018}. A função $h(x)$ é máxima quando
$e(x) = 0{,}5$ e tende a zero quando $e(x)$ se aproxima de 0 ou de 1. A população de
sobreposição dá, portanto, mais peso aos perfis em que os dois momentos de início são
frequentes e menos peso aos perfis em que um deles é quase certo.

Sob ignorabilidade condicional, positividade e consistência \cite{hernan2025},
$\tau_{\ATO}$ tem interpretação causal para essa população. Como essas hipóteses não
podem ser verificadas com os dados disponíveis, a estimativa é apresentada como diferença
de riscos ajustada: uma associação ajustada na população de sobreposição.

\subsection{Justificativa da escolha}

\citeonline{li2018} propuseram os pesos de sobreposição como membro de uma família de
pesos de balanceamento definida pela função $h$. Os autores mostram que, sob
homocedasticidade, esses pesos minimizam a variância assintótica do estimador ponderado
entre os pesos dessa família e que, com escore estimado por regressão logística, produzem
balanço exato das médias das covariáveis incluídas no modelo.

Neste estudo, três características do diagnóstico favoreceram o método. Primeiro, com
escores concentrados perto de 1, os pesos de probabilidade inversa para o ATE, proporcionais
a $1/(1-e)$ no grupo de comparação, ou para o ATT, proporcionais a $e/(1-e)$, seriam
elevados para os controles com escore alto. Segundo, recortar o suporte em limites
convencionais excluiria uma parcela grande do grupo $T=1$ e redefiniria a população-alvo
de forma arbitrária \cite{crump2009}. Terceiro, os pesos de sobreposição são limitados ao
intervalo $[0,1]$, não envolvem divisão por $e$ ou $1-e$ e reduzem gradualmente a
influência das regiões de pouca sobreposição, sem limiares escolhidos ad hoc. O custo
dessa escolha é que a população-alvo não é a população total, e precisa ser descrita
(Seção~\ref{sec:perfil}).

\subsection{Modelo de escore de propensão}

O escore foi estimado por regressão logística com as covariáveis da
Tabela~\ref{tab:variaveis}:
\begin{multline}\label{eq:logit}
\operatorname{logit} e(X_i) = \beta_0 + \beta_1\,\text{idade}_i + \beta_2\,\text{idade}_i^2
 + \sum_{k} \gamma_k\,\mathbb{1}\{\text{escolaridade}_i = k\}
 + \sum_{k} \delta_k\,\mathbb{1}\{\text{raça/cor}_i = k\} \\
 + \theta\,\mathbb{1}\{\text{nulípara}_i\}
 + \sum_{k} \lambda_k\,\mathbb{1}\{\text{perdas fetais}_i = k\},
\end{multline}
em que os indicadores incluem a categoria ignorado de cada variável. A estimação usou
penalização L2 (parâmetro $C = 1$ na parametrização do \emph{scikit-learn}), variáveis
numéricas padronizadas, intercepto não penalizado e nenhum peso de classe; com cerca de
115 mil registros, a penalização é praticamente desprezível. O modelo foi ajustado na
amostra inteira, e não fora da partição, porque a propriedade de balanço dos pesos de
sobreposição depende de o escore ser ajustado no mesmo conjunto em que os pesos são
aplicados. A otimização convergiu em 81 iterações, e a média do escore estimado (0,9118)
coincidiu com a proporção observada de $T=1$ (0,9118). O desfecho não foi usado na
especificação nem no ajuste.

\subsection{Pesos, riscos ponderados e diferença ajustada}

Com o escore estimado $\hat{e}(X_i)$, os pesos de sobreposição são
\begin{equation}\label{eq:pesos}
w_i =
\begin{cases}
1 - \hat{e}(X_i), & \text{se } T_i = 1,\\[2pt]
\hat{e}(X_i), & \text{se } T_i = 0.
\end{cases}
\end{equation}
Cada registro do grupo $T=1$ recebe peso igual à probabilidade estimada de pertencer ao
grupo de comparação, e cada registro de $T=0$, peso igual à probabilidade estimada de
pertencer ao grupo de início precoce. O risco ponderado em cada grupo e a diferença
ajustada são
\begin{equation}\label{eq:risco}
\hat{r}_t = \frac{\sum_{i:\,T_i = t} w_i\,Y_i}{\sum_{i:\,T_i = t} w_i},
\quad t \in \{0, 1\},
\qquad\qquad
\hat{\Delta}_{\ATO} = \hat{r}_1 - \hat{r}_0 ,
\end{equation}
com $\hat{\Delta}_{\ATO}$ expresso em pontos percentuais. Como os pesos são normalizados
dentro de cada grupo, $\hat{r}_t$ é a proporção ponderada de baixo peso no grupo $t$. Os
riscos brutos correspondem a $w_i = 1$ para todos os registros.

\subsection{Balanço ponderado e tamanho amostral efetivo}

O balanço após a ponderação foi avaliado pelas expressões \eqref{eq:smd-num} e
\eqref{eq:smd-cat}, com médias e proporções ponderadas no numerador e o mesmo denominador
não ponderado, para que a comparação antes/depois reflita apenas a mudança nas médias. A
perda de precisão decorrente da ponderação foi resumida pelo tamanho amostral efetivo de
Kish em cada grupo:
\begin{equation}\label{eq:kish}
n_{\mathrm{ef},t} = \frac{\left(\sum_{i:\,T_i = t} w_i\right)^2}{\sum_{i:\,T_i = t} w_i^2}.
\end{equation}

\subsection{Suporte com as covariáveis finais}

Como o modelo final difere do modelo de diagnóstico (sai a situação conjugal; entra o
termo quadrático de idade), o suporte foi reexaminado (Tabela~\ref{tab:suporte}). A regra
foi fixada antes do cálculo de qualquer risco: nenhum recorte automático; contar os
registros fora do suporte comum observado e sua participação na soma dos pesos. Seis
registros de $T=1$ ficaram abaixo do mínimo dos controles, com 0,022\% da soma dos pesos do
grupo. Nenhum registro foi excluído.

\begin{table}[htbp]
\centering
\caption{Distribuição do escore de propensão estimado no modelo final, por grupo}
\label{tab:suporte}
\small
\begin{tabular}{@{}l *{9}{S[table-format=1.4]}@{}}
\toprule
\cab{Grupo} & {\cab{Mín.}} & {\cab{P1}} & {\cab{P5}} & {\cab{P25}} & {\cab{Mediana}} &
  {\cab{P75}} & {\cab{P95}} & {\cab{P99}} & {\cab{Máx.}} \\
\midrule
$T=1$ & 0.6232 & 0.8027 & 0.8450 & 0.8870 & 0.9124 & 0.9530 & 0.9719 & 0.9725 & 0.9760 \\
$T=0$ & 0.6910 & 0.7732 & 0.8192 & 0.8678 & 0.8936 & 0.9169 & 0.9643 & 0.9719 & 0.9760 \\
\bottomrule
\end{tabular}
\nota{Suporte comum observado: 0,6910 a 0,9760. Fora do suporte: 6 registros de $T=1$ e
nenhum de $T=0$.}
\fonte{\arquivo{notebooks/principal/06_ato_sp_capital_2024.ipynb} (seção 2).}
\end{table}

\subsection{Propriedades e limites do estimador}\label{sec:propriedades}

Duas propriedades precisam constar do texto. Primeiro, com escore logístico, as médias
ponderadas dos termos incluídos no modelo tornam-se iguais entre os grupos \cite{li2018};
com a penalização L2, aproximadamente iguais. O balanço quase exato nas covariáveis
ajustadas é, portanto, esperado por construção e não demonstra ausência de confundimento:
o diagnóstico informativo está nas covariáveis não incluídas e no tamanho amostral
efetivo. Segundo, com um único modelo de escore e sem modelo para o desfecho, o estimador
\emph{não é duplamente robusto}. Sua validade depende de o escore logístico de efeitos
principais, com idade quadrática, representar adequadamente a relação entre $X$ e $T$.
Estimadores duplamente robustos, como o AIPW usado na fase anterior do projeto, combinam
modelos de escore e de desfecho \cite{funk2011}, mas também não eliminam confundimento
não medido.

% =====================================================================
\section{Incerteza estatística}\label{sec:incerteza}

A incerteza foi quantificada por bootstrap não paramétrico com $B = 1.000$ réplicas. Em
cada réplica $b$:
\begin{enumerate}
\item sorteiam-se, com reposição, 114.694 registros da população analítica, o mesmo
tamanho da base original;
\item o escore de propensão é reajustado na réplica, com a mesma especificação do modelo
principal;
\item os pesos, os riscos e as diferenças são recalculados, produzindo
$\hat{\Delta}^{*}_{b}$.
\end{enumerate}
O intervalo de confiança de 95\% é formado pelos percentis 2,5 e 97,5 da distribuição das
réplicas,
\begin{equation}\label{eq:ic}
\mathrm{IC}_{95\%} = \left[\, \hat{\Delta}^{*}_{(0{,}025)}\,,\; \hat{\Delta}^{*}_{(0{,}975)} \,\right].
\end{equation}
Como conferência, calcularam-se o erro-padrão bootstrap e o intervalo normal
correspondente. As sementes aleatórias foram derivadas de uma sequência fixa, uma por
réplica, de modo que o resultado não depende do número de processos paralelos; todas as
1.000 réplicas convergiram. Reajustar o escore em cada réplica incorpora ao intervalo a
variabilidade decorrente da estimação dos pesos.

O bootstrap não deve ser confundido com amostragem dos dados. A análise principal usou
todos os 114.694 registros elegíveis; não houve sorteio de uma subamostra, por exemplo de
10 mil nascimentos, para estimar o resultado. A reamostragem bootstrap é um procedimento
computacional aplicado depois da estimação: cada réplica tem o mesmo tamanho da base,
repete alguns registros e omite outros, e serve apenas para aproximar a variabilidade do
estimador. As réplicas não são novas amostras da população nem substituem a estimativa
pontual, que é sempre a obtida com a base completa. Como a base reúne todos os registros
elegíveis do ano, o intervalo deve ser lido como medida da variabilidade do estimador sob
um modelo de nascimentos independentes, e não como erro de amostragem de um levantamento.

Essa independência é uma simplificação. Agrupamentos por estabelecimento de saúde, por
região de residência ou por mães com mais de um parto no ano tendem a aumentar a
incerteza, e ignorá-los pode produzir intervalos estreitos demais \cite{cameron2015}.
Além disso, o intervalo reflete apenas variabilidade amostral; não incorpora vieses de
confundimento, mensuração ou seleção.

O procedimento bootstrap foi aplicado somente à estimação principal de 2024. A estimativa
de sensibilidade (Seção~\ref{sec:sensibilidade}) e as estimativas anuais de 2014 a 2024
(Seção~\ref{sec:serie}) são pontuais, sem intervalo de confiança.

% =====================================================================
\section{Resultados}\label{sec:resultados}

\subsection{População analisada}

A população analítica reuniu 114.694 nascimentos, 104.576 (91,18\%) no grupo $T=1$ e
10.118 (8,82\%) no grupo $T=0$, com 9.850 casos de baixo peso ao nascer
(Tabela~\ref{tab:populacao}). Nenhum registro foi excluído por suporte ou por covariável
incompleta.

\begin{table}[htbp]
\centering
\caption{População analisada e casos de baixo peso ao nascer, São Paulo (capital), 2024}
\label{tab:populacao}
\small
\begin{tabular}{@{}l S[table-format=6.0] S[table-format=3.2] S[table-format=4.0]@{}}
\toprule
\cab{Grupo} & {\cab{Nascimentos}} & {\cab{\% do total}} & {\cab{Casos de baixo peso}} \\
\midrule
$T=1$ (início nos meses 1--3) & 104576 & 91.18 & 8797 \\
$T=0$ (início nos meses 4--6) & 10118  & 8.82  & 1053 \\
\midrule
Total                         & 114694 & 100.00 & 9850 \\
\bottomrule
\end{tabular}
\fonte{\arquivo{notebooks/principal/06_ato_sp_capital_2024.ipynb} (seções 1 e 5).}
\end{table}

\subsection{Balanço das covariáveis}

Nas covariáveis incluídas no modelo, o maior $|\mathrm{SMD}|$ caiu de 0,452 para 0,0031
(Figura~\ref{fig:balanco}), resultado esperado por construção. O indicador descritivo de
mães de 10--19 anos, que não entra no modelo, também ficou equilibrado: 11,1\% em $T=1$ e
11,3\% em $T=0$ após a ponderação. A situação conjugal, deixada fora do modelo principal,
permaneceu desequilibrada: o SMD do nível casada passou de 0,421 para 0,194; o de
solteira, de $-0{,}279$ para $-0{,}117$; e o de união estável, de $-0{,}125$ para
$-0{,}070$. A ponderação por covariáveis correlacionadas reduziu, mas não eliminou, o
desequilíbrio em uma característica não ajustada. Entre as linhas não ajustadas ou
descritivas, o maior $|\mathrm{SMD}|$ após a ponderação foi 0,194.

\begin{figure}[htbp]
\centering
\caption{Diferença média padronizada absoluta por nível de covariável, antes e depois da
ponderação ATO}
\label{fig:balanco}
% Figura 2 — |SMD| por nível de covariável, antes e depois da ponderação ATO.
% Valores copiados do gráfico Plotly salvo em
% notebooks/principal/06_ato_sp_capital_2024.ipynb (seção 4).
% SMD com denominador não ponderado, igual antes e depois.
% Colunas: linha (posição vertical, 34 = topo), antes, depois.
\pgfplotstableread{
linha antes   depois
33    0.2703  0.0002
32    0.2529  0.0001
31    0.1857  0.0064
29    0.0116  0.0003
28    0.0619  0.0003
27    0.2057  0.0001
26    0.2338  0.0002
25    0.0137  0.0003
24    0.4518  0.0003
23    0.0215  0.0031
21    0.1990  0.0000
20    0.0558  0.0002
19    0.0567  0.0006
18    0.1678  0.0002
17    0.0322  0.0013
16    0.0042  0.0013
14    0.1659  0.0002
13    0.1659  0.0002
11    0.0375  0.0001
10    0.0414  0.0002
9     0.0002  0.0000
8     0.0035  0.0003
6     0.2785  0.1172
5     0.4214  0.1945
4     0.0001  0.0020
3     0.0218  0.0031
2     0.1254  0.0704
1     0.0101  0.0090
}\dadosbalanco
\newcommand{\grupo}[1]{\textbf{\color{iptazul}#1}}
\newcommand{\nivel}[1]{\hspace*{0.9em}#1}
\newcommand{\seg}[3]{\draw[iptcinza!45, line width=0.7pt] (axis cs:#2,#1) -- (axis cs:#3,#1);}
\begin{tikzpicture}
\begin{axis}[
  estilo ipt,
  set layers,
  scale only axis,
  width=7.9cm, height=12.0cm,
  xmin=-0.012, xmax=0.48,
  ymin=0.4, ymax=34.6,
  xtick={0,0.1,0.2,0.3,0.4},
  xticklabel style={/pgf/number format/.cd, fixed, precision=1, use comma},
  xmajorgrids,
  ytick={1,...,34},
  major y tick style={draw=none},
  yticklabel style={text width=5.3cm, align=left, font=\sffamily\footnotesize},
  yticklabels={
    \nivel{Ignorado/ausente}, \nivel{União estável}, \nivel{Separada/divorciada},
    \nivel{Viúva}, \nivel{Casada}, \nivel{Solteira},
    \grupo{Situação conjugal (não ajustada)},
    \nivel{Ignorado/ausente}, \nivel{Duas ou mais}, \nivel{Uma}, \nivel{Nenhuma},
    \grupo{Perdas fetais anteriores},
    \nivel{Multípara}, \nivel{Nulípara},
    \grupo{Paridade},
    \nivel{Ignorado/ausente}, \nivel{Indígena}, \nivel{Parda}, \nivel{Amarela},
    \nivel{Preta}, \nivel{Branca},
    \grupo{Raça/cor},
    \nivel{Ignorado/ausente}, \nivel{Superior completo}, \nivel{Superior incompleto},
    \nivel{Médio}, \nivel{Fundamental II}, \nivel{Fundamental I}, \nivel{Sem escolaridade},
    \grupo{Escolaridade},
    \nivel{10–19 anos (descritiva)}, \nivel{Quadrado da idade (média)}, \nivel{Média (anos)},
    \grupo{Idade materna},
  },
  xlabel={Diferença média padronizada absoluta (|SMD|)},
  legend style={at={(0.5,1.01)}, anchor=south, legend columns=2,
    /tikz/every even column/.append style={column sep=1.2em}},
]
% Fundo da covariável não incluída no modelo principal
\begin{pgfonlayer}{axis background}
  \fill[iptazul!6] (axis cs:-0.012,0.4) rectangle (axis cs:0.48,7.5);
  % separadores entre grupos de covariáveis
  \draw[iptcinza!35, line width=0.4pt] (axis cs:-0.012,7.5) -- (axis cs:0.48,7.5);
  \draw[iptcinza!35, line width=0.4pt] (axis cs:-0.012,12.5) -- (axis cs:0.48,12.5);
  \draw[iptcinza!35, line width=0.4pt] (axis cs:-0.012,15.5) -- (axis cs:0.48,15.5);
  \draw[iptcinza!35, line width=0.4pt] (axis cs:-0.012,22.5) -- (axis cs:0.48,22.5);
  \draw[iptcinza!35, line width=0.4pt] (axis cs:-0.012,30.5) -- (axis cs:0.48,30.5);
  \draw[iptcinza!80, dashed, line width=0.6pt] (axis cs:0.1,0.4) -- (axis cs:0.1,34.6);
\end{pgfonlayer}
% Segmentos ligando antes e depois: \seg{linha}{antes}{depois}
\seg{33}{0.2703}{0.0002} \seg{32}{0.2529}{0.0001} \seg{31}{0.1857}{0.0064}
\seg{29}{0.0116}{0.0003} \seg{28}{0.0619}{0.0003} \seg{27}{0.2057}{0.0001}
\seg{26}{0.2338}{0.0002} \seg{25}{0.0137}{0.0003} \seg{24}{0.4518}{0.0003}
\seg{23}{0.0215}{0.0031} \seg{21}{0.1990}{0.0000} \seg{20}{0.0558}{0.0002}
\seg{19}{0.0567}{0.0006} \seg{18}{0.1678}{0.0002} \seg{17}{0.0322}{0.0013}
\seg{16}{0.0042}{0.0013} \seg{14}{0.1659}{0.0002} \seg{13}{0.1659}{0.0002}
\seg{11}{0.0375}{0.0001} \seg{10}{0.0414}{0.0002} \seg{9}{0.0002}{0.0000}
\seg{8}{0.0035}{0.0003}  \seg{6}{0.2785}{0.1172}  \seg{5}{0.4214}{0.1945}
\seg{4}{0.0001}{0.0020}  \seg{3}{0.0218}{0.0031}  \seg{2}{0.1254}{0.0704}
\seg{1}{0.0101}{0.0090}
\addplot[only marks, mark=o, mark size=2.4pt, cinzaserie, line width=0.9pt]
  table[x=antes, y=linha] {\dadosbalanco};
\addlegendentry{Antes da ponderação}
\addplot[only marks, mark=*, mark size=2.2pt, iptazul]
  table[x=depois, y=linha] {\dadosbalanco};
\addlegendentry{Após ponderação ATO}
\end{axis}
\end{tikzpicture}

\nota{Linha tracejada: referência convencional de 0,1. Faixa sombreada: situação conjugal,
não incluída no modelo principal. Denominador do SMD não ponderado, igual antes e depois.}
\fonte{\arquivo{notebooks/principal/06_ato_sp_capital_2024.ipynb} (seção 4).}
\end{figure}

\subsection{Pesos e tamanho amostral efetivo}

Os pesos ficaram concentrados em faixas estreitas e não houve valores extremos
(Tabela~\ref{tab:pesos}). O tamanho amostral efetivo foi de 84.995 em $T=1$ (81,3\% do
tamanho do grupo) e de 10.096 em $T=0$ (99,8\%). A maior participação de um registro
individual na soma dos pesos do grupo foi de 0,0042\% em $T=1$ e de 0,0108\% em $T=0$. A
ponderação implicou perda moderada de precisão em $T=1$ e praticamente nenhuma em $T=0$.

\begin{table}[htbp]
\centering
\caption{Distribuição dos pesos de sobreposição e tamanho amostral efetivo, por grupo}
\label{tab:pesos}
\small
\begin{tabular}{@{}l c c@{}}
\toprule
\cab{Medida} & \cab{$T=1$ (peso $1-\hat{e}$)} & \cab{$T=0$ (peso $\hat{e}$)} \\
\midrule
Registros ($n$)                               & 104.576         & 10.118 \\
Tamanho efetivo de Kish ($n_{\mathrm{ef}}$)   & 84.995          & 10.096 \\
$n_{\mathrm{ef}}/n$                           & 81,3\%          & 99,8\% \\
Peso mínimo -- máximo                         & 0,0240 -- 0,3768 & 0,6910 -- 0,9760 \\
Peso mediano                                  & 0,0876          & 0,8936 \\
Maior participação individual na soma dos pesos & 0,0042\%      & 0,0108\% \\
\bottomrule
\end{tabular}
\fonte{\arquivo{notebooks/principal/06_ato_sp_capital_2024.ipynb} (seção 3).}
\end{table}

\subsection{Perfil da população de sobreposição}\label{sec:perfil}

Como o balanço é praticamente exato nas covariáveis ajustadas, as médias ponderadas de
$T=1$ e $T=0$ coincidem e descrevem a população de sobreposição. Comparada à população
total, ela é mais jovem e menos escolarizada, com perfil próximo ao do grupo $T=0$
(Tabela~\ref{tab:perfil}). O resultado apresentado a seguir vale para essa população: idade média de
27,6 anos e 13,8\% com ensino superior completo, contra 29,1 anos e 29,5\% na população
total.

\begin{table}[htbp]
\centering
\caption{Perfil da população total, dos grupos sem ponderação e da população de
sobreposição}
\label{tab:perfil}
\small
\begin{tabular}{@{}l S[table-format=2.1] S[table-format=2.1] S[table-format=2.1]
  S[table-format=2.1]@{}}
\toprule
\cab{Característica} & {\cab{População total}} & {\cab{$T=1$ bruto}} & {\cab{$T=0$ bruto}}
  & {\cab{Sobreposição (ATO)}} \\
\midrule
Idade materna média (anos)         & 29.1 & 29.3 & 27.5 & 27.6 \\
Mães de 10--19 anos (\%)           & 7.1  & 6.6  & 11.9 & 11.1 \\
Escolaridade: ensino médio (\%)    & 55.0 & 54.0 & 65.4 & 65.2 \\
Escolaridade: superior completo (\%) & 29.5 & 31.1 & 12.9 & 13.8 \\
Raça/cor branca (\%)               & 45.0 & 45.9 & 36.1 & 36.7 \\
Raça/cor parda (\%)                & 42.7 & 41.9 & 50.3 & 49.7 \\
Nulíparas (\%)                     & 39.7 & 40.4 & 32.4 & 32.9 \\
\bottomrule
\end{tabular}
\nota{Coluna de sobreposição: médias ponderadas de $T=1$, iguais às de $T=0$ nas
covariáveis ajustadas; para mães de 10--19 anos, indicador não ajustado, o valor de $T=0$
ponderado é 11,3\%.}
\fonte{\arquivo{notebooks/principal/06_ato_sp_capital_2024.ipynb} (seção 4).}
\end{table}

\subsection{Riscos brutos e ponderados}

O risco bruto de baixo peso foi 8,41\% em $T=1$ e 10,41\% em $T=0$, diferença de
$-2{,}00$\pp{} (IC 95\%: $-2{,}61$ a $-1{,}37$). Após a ponderação ATO, os riscos foram
8,68\% e 10,42\%, e a diferença ajustada foi de $-1{,}74$\pp{} (IC 95\%: $-2{,}36$ a
$-1{,}12$), conforme a Tabela~\ref{tab:riscos} e a Figura~\ref{fig:riscos}. O erro-padrão
bootstrap da diferença ajustada foi 0,32\pp, o intervalo normal de conferência foi de
$-2{,}37$ a $-1{,}11$\pp{} e a média das réplicas, $-1{,}75$\pp{} O intervalo não inclui
zero. A ponderação reduziu a magnitude da diferença bruta em cerca de 13\%.

\begin{table}[htbp]
\centering
\caption{Riscos de baixo peso ao nascer, brutos e ponderados (ATO), e diferença entre
grupos, São Paulo (capital), 2024}
\label{tab:riscos}
\small
\setlength{\tabcolsep}{5pt}
\begin{tabular}{@{}l S[table-format=2.2] c S[table-format=2.2] c S[table-format=-1.2] c@{}}
\toprule
 & \multicolumn{2}{c}{\cab{$T=1$ (meses 1--3)}} & \multicolumn{2}{c}{\cab{$T=0$ (meses 4--6)}}
 & \multicolumn{2}{c}{\cab{Diferença $T=1 - T=0$}} \\
\cmidrule(lr){2-3}\cmidrule(lr){4-5}\cmidrule(l){6-7}
\cab{Estimativa} & {\cab{Risco (\%)}} & \cab{IC 95\%} & {\cab{Risco (\%)}} & \cab{IC 95\%}
 & {\cab{p.p.}} & \cab{IC 95\%} \\
\midrule
Bruta              & 8.41 & 8,25 a 8,59 & 10.41 & 9,84 a 11,00 & -2.00 & $-$2,61 a $-$1,37 \\
\destaque
Ponderada (ATO)    & 8.68 & 8,49 a 8,87 & 10.42 & 9,85 a 11,00 & -1.74 & $-$2,36 a $-$1,12 \\
\bottomrule
\end{tabular}
\nota{Intervalos bootstrap por percentis, 1.000 réplicas com reestimação do escore. Riscos
brutos sobre todos os nascimentos de cada grupo; riscos ATO com pesos normalizados dentro
do grupo. Diferença ATO: erro-padrão bootstrap 0,32\pp; intervalo normal $-2{,}37$ a
$-1{,}11$\pp}
\fonte{\arquivo{notebooks/principal/06_ato_sp_capital_2024.ipynb} (seção 5).}
\end{table}

\begin{figure}[htbp]
\centering
\caption{Risco de baixo peso ao nascer por grupo, bruto e ponderado (ATO), e diferença
entre grupos, com intervalos de confiança de 95\%}
\label{fig:riscos}
% Figura 3 — Riscos brutos e ponderados (ATO) e diferenças T=1 − T=0, com IC 95%
% bootstrap por percentis (1.000 réplicas, escore reestimado em cada réplica).
% Painel (a): valores do gráfico Plotly salvo em
%   notebooks/principal/06_ato_sp_capital_2024.ipynb (seção 5), em pontos percentuais.
% Painel (b): diferenças e intervalos da tabela de riscos do mesmo notebook (seção 5).
% Os dois painéis são alinhados pela base da área de plotagem (baseline=(eixo.south)).
\begin{minipage}[b]{0.49\linewidth}
\centering
\begin{tikzpicture}[baseline=(eixo.south)]
\begin{axis}[
  name=eixo,
  estilo ipt,
  scale only axis,
  width=5.0cm, height=4.8cm,
  title={(a) Risco de baixo peso por grupo},
  xmin=0.45, xmax=2.55,
  xtick={1,2},
  xticklabels={{$T=1$\\meses 1–3}, {$T=0$\\meses 4–6}},
  xticklabel style={align=center},
  ymin=7.6, ymax=11.4,
  ytick={8,9,10,11},
  ymajorgrids,
  ylabel={Risco de peso $<$ 2.500 g (\%)},
  legend style={at={(0.03,0.97)}, anchor=north west},
  error bars/error bar style={line width=0.9pt},
  error bars/error mark options={rotate=90, mark size=2.6pt, line width=0.9pt},
]
\addplot[serie secundaria, only marks, mark=o, mark size=2.6pt,
  error bars/.cd, y dir=both, y explicit]
  coordinates {(0.86,8.4121) += (0,0.1777) -= (0,0.1647)
               (1.86,10.4072) += (0,0.5903) -= (0,0.5698)};
\addlegendentry{Bruto}
\addplot[serie principal, only marks, mark=diamond*, mark size=3.2pt,
  error bars/.cd, y dir=both, y explicit]
  coordinates {(1.14,8.6809) += (0,0.1853) -= (0,0.1916)
               (2.14,10.4183) += (0,0.5835) -= (0,0.5720)};
\addlegendentry{Ponderado (ATO)}
\node[rotulo valor, anchor=east] at (axis cs:0.80,8.4121) {8,41};
\node[rotulo valor, anchor=west] at (axis cs:1.20,8.6809) {8,68};
\node[rotulo valor, anchor=east] at (axis cs:1.80,10.4072) {10,41};
\node[rotulo valor, anchor=west] at (axis cs:2.20,10.4183) {10,42};
\end{axis}
\end{tikzpicture}
\end{minipage}%
\hfill
\begin{minipage}[b]{0.49\linewidth}
\centering
\begin{tikzpicture}[baseline=(eixo.south)]
\begin{axis}[
  name=eixo,
  estilo ipt,
  scale only axis,
  width=5.0cm, height=4.8cm,
  title={(b) Diferença $T=1$ − $T=0$ (p.p.)},
  xmin=-3.2, xmax=0.6,
  xtick={-3,-2,-1,0},
  xmajorgrids,
  ymin=0.35, ymax=2.75,
  ytick={1,2},
  yticklabels={Ponderada (ATO), Bruta},
  yticklabel style={font=\sffamily\footnotesize},
  major y tick style={draw=none},
  xlabel={Diferença de riscos (p.p.)},
  error bars/error bar style={line width=0.9pt},
  error bars/error mark options={rotate=90, mark size=2.6pt, line width=0.9pt},
]
\draw[iptcinza, dashed, line width=0.6pt] (axis cs:0,0.35) -- (axis cs:0,2.75);
\addplot[serie secundaria, only marks, mark=o, mark size=2.6pt,
  error bars/.cd, x dir=both, x explicit]
  coordinates {(-2.00,2) += (0.63,0) -= (0.61,0)};
\addplot[serie principal, only marks, mark=diamond*, mark size=3.2pt,
  error bars/.cd, x dir=both, x explicit]
  coordinates {(-1.74,1) += (0.62,0) -= (0.62,0)};
% Rótulos centrados na área de plotagem, acima de cada estimativa
\node[rotulo valor, anchor=south] at (axis cs:-1.30,2.16)
  {$-$2,00 \ (IC 95\%: $-$2,61 a $-$1,37)};
\node[rotulo valor, anchor=south, font=\sffamily\scriptsize\bfseries, text=iptazul]
  at (axis cs:-1.30,1.16) {$-$1,74 \ (IC 95\%: $-$2,36 a $-$1,12)};
\end{axis}
\end{tikzpicture}
\end{minipage}

\nota{Barras: IC 95\% bootstrap por percentis (1.000 réplicas, escore reestimado em cada
réplica). Linha tracejada no painel (b): diferença nula. População: 114.694 nascimentos de
mães residentes no município de São Paulo, 2024.}
\fonte{\arquivo{notebooks/principal/06_ato_sp_capital_2024.ipynb} (seção 5).}
\end{figure}

A mudança concentrou-se no grupo $T=1$, cujo risco ponderado (8,68\%) é maior que o
bruto (8,41\%). Esse movimento acompanha a mudança de composição do grupo, que após a
ponderação passa a refletir o perfil da população de sobreposição, mais jovem e menos
escolarizada (Seção~\ref{sec:perfil}). No grupo $T=0$, em que os pesos são próximos de
uniformes, o risco praticamente não se altera (10,41\% e 10,42\%).

Em síntese: \emph{na população de
sobreposição definida pelos pesos ATO, após o equilíbrio das distribuições de idade,
escolaridade, raça/cor, paridade e perdas fetais anteriores, o risco de baixo peso ao
nascer foi 1,74 p.p.\ menor entre nascimentos com início do pré-natal nos meses 1--3 do
que entre nascimentos com início nos meses 4--6 (IC 95\%: $-2{,}36$ a $-1{,}12$ p.p.).}
Trata-se de associação ajustada por características observadas, não de efeito causal
identificado.

% =====================================================================
\section{Análise de sensibilidade: situação conjugal}\label{sec:sensibilidade}

A situação conjugal foi o segundo maior desequilíbrio observado e permaneceu
desequilibrada no modelo principal. Para avaliar quanto a estimativa depende de sua
exclusão, a variável foi acrescentada ao escore em uma análise exploratória, mantidos a
população, as codificações, a idade e seu quadrado, a regressão logística com penalização
L2, o ajuste na amostra inteira, os pesos de sobreposição e a ausência de recorte de
suporte. O novo escore convergiu em 89 iterações.

\begin{table}[htbp]
\centering
\caption{Modelo principal e modelo de sensibilidade com situação conjugal no escore}
\label{tab:sensibilidade}
\small
\begin{tabular}{@{}l c c@{}}
\toprule
\cab{Medida} & \cab{Principal} & \cab{Com situação conjugal} \\
\midrule
Diferença $T=1 - T=0$ (p.p.)                  & $-$1,74 (IC 95\%: $-$2,36 a $-$1,12) & $-$1,63 (sem IC) \\
Risco ponderado $T=1$ / $T=0$ (\%)            & 8,68 / 10,42   & 8,77 / 10,40 \\
Maior $|\mathrm{SMD}|$, conjunto diagnóstico comum & 0,1945    & 0,0111 \\
SMD do nível casada                           & 0,1945         & 0,0001 \\
Tamanho efetivo $T=1$ / $T=0$                 & 84.995 / 10.096 & 82.323 / 10.094 \\
\bottomrule
\end{tabular}
\nota{Conjunto diagnóstico comum: mesmas linhas nos dois modelos, incluindo situação
conjugal e o indicador de 10--19 anos. Na sensibilidade, o maior desequilíbrio residual foi
o desse indicador ($-0{,}0111$).}
\fonte{\arquivo{notebooks/principal/06_ato_sp_capital_2024.ipynb} (seção 7).}
\end{table}

A diferença passou de $-1{,}74$ para $-1{,}63$\pp, uma redução de 0,107\pp{} (6,1\%) na
magnitude absoluta, na mesma direção (Tabela~\ref{tab:sensibilidade}). O SMD do nível
casada aproximou-se de zero, como esperado ao incluir a variável no modelo. Essa
comparação exige quatro cautelas.

\begin{itemize}
\item \emph{Ausência de intervalo de confiança.} Não foi executado bootstrap para a
sensibilidade, e o intervalo do modelo principal não se aplica à nova estimativa. A
comparação é apenas entre estimativas pontuais, sem teste da diferença entre modelos; a
semelhança não constitui teste de equivalência ou de robustez causal.
\item \emph{Temporalidade.} A situação conjugal é registrada no parto, e não foi
demonstrado que ela anteceda o início do pré-natal.
\item \emph{Mudança da população-alvo.} Incluir a variável altera o escore e, portanto, os
pesos e a população de sobreposição, cuja densidade é proporcional a
$f(x)\,e(x)\,[1-e(x)]$. A proporção de casadas na população ponderada passou a 19,2\% nos
dois grupos, enquanto no modelo principal era de 27,1\% em $T=1$ e 18,6\% em $T=0$.
\item \emph{Atribuição.} A diferença entre as estimativas mistura mudança de ajuste e
mudança de população e não pode ser atribuída exclusivamente à correção de confundimento.
\end{itemize}

% =====================================================================
\section{Evolução histórica da associação ajustada, 2014--2024}\label{sec:serie}

\subsection{Especificação anual}

Para descrever se a direção e a magnitude da associação ajustada se mantêm ao longo do
período, a especificação principal foi repetida separadamente em cada ano de 2014 a 2024:
os mesmos critérios de população, exposição e desfecho; as mesmas covariáveis e
codificações; o mesmo escore logístico com penalização L2, ajustado na amostra inteira de
cada ano; os mesmos pesos de sobreposição; e nenhum recorte de suporte. Cada ano tem seu
próprio modelo de escore, e não há modelo combinado entre anos. A estimativa de 2024
reproduz a da Seção~\ref{sec:resultados}. Os onze ajustes convergiram, com 53 a 97
iterações e sem avisos.

Em três anos, um registro com idade materna fora do intervalo de 10 a 59 anos --- 60
anos em 2017, 62 em 2022 e 61 em 2023 --- foi mantido. Para esses registros, idade e
idade$^2$ receberam a imputação pela mediana anual e o indicador de ausência previstos no
procedimento de estimação do escore, mediante validação manual. Nenhum registro foi
excluído. As regras operacionais de codificação foram preservadas em todos os anos, sem
que se tenha comprovado a invariância histórica do preenchimento dos campos.

A série é descritiva. Não foi executado bootstrap, não foram calculados intervalos de
confiança anuais e não há teste de diferença entre anos. O intervalo de $-2{,}36$ a
$-1{,}12$\pp{} pertence apenas à estimação principal de 2024 e não deve ser estendido aos
demais anos.

\subsection{Estimativas anuais}

A diferença ajustada foi negativa nos onze anos (Tabela~\ref{tab:serie} e
Figura~\ref{fig:serie}). A magnitude variou de $-0{,}58$\pp{} em 2020 a $-1{,}90$\pp{} em
2023, amplitude de 1,33\pp, sem tendência monotônica: 2014--2016 ficaram próximos de
$-1$\pp; 2017--2018 tiveram magnitudes maiores; 2019--2020 aproximaram-se de zero; e
2021--2024 voltaram a valores mais negativos. Em 2024, a estimativa de $-1{,}74$\pp{}
ficou próxima das de 2021 ($-1{,}73$) e 2022 ($-1{,}77$). Em todos os anos, a diferença
ponderada teve magnitude menor que a bruta, como em 2024.

\begin{table}[htbp]
\centering
\caption{Estimativas anuais com a especificação principal, São Paulo (capital), 2014--2024}
\label{tab:serie}
\footnotesize
\setlength{\tabcolsep}{3.6pt}
\begin{tabular}{@{}l S[table-format=6.0] S[table-format=5.0] S[table-format=2.2]
  S[table-format=-1.2] S[table-format=1.2] S[table-format=2.2] S[table-format=-1.2]
  S[table-format=6.0] S[table-format=5.0]@{}}
\toprule
 & \multicolumn{3}{c}{\cab{População}} & {\cab{Bruta}} &
  \multicolumn{3}{c}{\cab{Ponderada (ATO)}} & \multicolumn{2}{c}{\cab{$n_{\mathrm{ef}}$ de Kish}} \\
\cmidrule(lr){2-4}\cmidrule(lr){5-5}\cmidrule(lr){6-8}\cmidrule(l){9-10}
\cab{Ano} & {\cab{Elegíveis}} & {\cab{$T=0$}} & {\cab{\% $T=0$}} & {\cab{Dif.}} &
  {\cab{Risco $T=1$}} & {\cab{Risco $T=0$}} & {\cab{Dif.}} & {\cab{$T=1$}} & {\cab{$T=0$}} \\
\midrule
2014 & 164987 & 24348 & 14.76 & -1.28 & 7.61 & 8.72  & -1.11 & 108879 & 24155 \\
2015 & 165787 & 24345 & 14.68 & -1.17 & 7.61 & 8.60  & -0.99 & 110184 & 24153 \\
2016 & 158025 & 21102 & 13.35 & -1.32 & 7.80 & 8.88  & -1.08 & 108773 & 20978 \\
2017 & 160429 & 18959 & 11.82 & -1.77 & 7.68 & 9.24  & -1.57 & 112037 & 18879 \\
2018 & 156806 & 17814 & 11.36 & -2.05 & 7.75 & 9.55  & -1.80 & 108388 & 17733 \\
2019 & 150427 & 17581 & 11.69 & -1.43 & 7.87 & 9.06  & -1.19 & 106088 & 17509 \\
2020 & 139283 & 15443 & 11.09 & -0.75 & 7.62 & 8.20  & -0.58 & 99307  & 15389 \\
2021 & 129761 & 12705 & 9.79  & -1.97 & 8.04 & 9.77  & -1.73 & 92163  & 12666 \\
2022 & 125306 & 12370 & 9.87  & -2.06 & 8.46 & 10.23 & -1.77 & 88856  & 12331 \\
2023 & 122244 & 11470 & 9.38  & -2.20 & 8.66 & 10.56 & -1.90 & 89016  & 11440 \\
\destaque
2024 & 114694 & 10118 & 8.82  & -2.00 & 8.68 & 10.42 & -1.74 & 84995  & 10096 \\
\bottomrule
\end{tabular}
\nota{Diferenças $T=1 - T=0$ em pontos percentuais; riscos em \%. \% $T=0$: sobre os
elegíveis do ano. Riscos brutos sobre todos os nascimentos do grupo; riscos ATO com pesos
normalizados dentro do grupo. Estimativas pontuais de onze modelos anuais independentes,
sem intervalo de confiança; o IC 95\% da Seção~\ref{sec:resultados} vale apenas para 2024.}
\fonte{\arquivo{notebooks/principal/07_evolucao_ato_sp_capital_2014_2024.ipynb} (seção 3).}
\end{table}

\begin{figure}[htbp]
\centering
\caption{Diferença de riscos de baixo peso ao nascer entre $T=1$ e $T=0$, bruta e
ponderada (ATO), e participação do grupo de comparação, por ano, São Paulo (capital),
2014--2024}
\label{fig:serie}
% Figura 4 (Etapa 2) — Série anual 2014–2024 da especificação principal.
% Painel (a): diferenças T=1 − T=0 (p.p.), ponderada (ATO) e bruta, copiadas do gráfico
%   Plotly salvo em notebooks/principal/07_evolucao_ato_sp_capital_2014_2024.ipynb (seção 5).
% Painel (b): % de T=0 entre os elegíveis do ano, do gráfico Plotly do mesmo notebook
%   (seção 6).
% Estimativas pontuais de onze modelos anuais independentes, sem intervalo de confiança.
\begin{tikzpicture}
\begin{axis}[
  estilo ipt,
  width=0.94\linewidth, height=6.0cm,
  title={(a) Diferença de riscos $T=1$ − $T=0$ (p.p.)},
  xmin=2013.6, xmax=2024.4,
  xtick={2014,...,2024},
  xticklabel style={/pgf/number format/1000 sep={}},
  ymin=-2.5, ymax=0.4,
  ytick={-2.5,-2,...,0},
  ymajorgrids,
  ylabel={Diferença de riscos (p.p.)},
  legend style={at={(0.02,0.97)}, anchor=north west},
]
\draw[iptcinza, dashed, line width=0.6pt] (axis cs:2013.6,0) -- (axis cs:2024.4,0);
\addplot[serie secundaria, dashed, mark=o] coordinates {
  (2014,-1.2846) (2015,-1.1740) (2016,-1.3195) (2017,-1.7729) (2018,-2.0461)
  (2019,-1.4339) (2020,-0.7471) (2021,-1.9731) (2022,-2.0564) (2023,-2.1965)
  (2024,-1.9951)};
\addlegendentry{Bruta (sem ponderação)}
\addplot[serie principal, mark=*] coordinates {
  (2014,-1.1069) (2015,-0.9870) (2016,-1.0788) (2017,-1.5665) (2018,-1.8002)
  (2019,-1.1886) (2020,-0.5776) (2021,-1.7269) (2022,-1.7690) (2023,-1.9027)
  (2024,-1.7374)};
\addlegendentry{Ponderada (ATO)}
% Rótulos seletivos da série ATO: primeiro ano, extremos e último ano
\node[rotulo valor, anchor=south, yshift=2pt] at (axis cs:2014,-1.1069) {$-$1,11};
\node[rotulo valor, anchor=south, yshift=2pt] at (axis cs:2020,-0.5776) {$-$0,58};
\node[rotulo valor, anchor=south, yshift=2pt] at (axis cs:2023,-1.9027) {$-$1,90};
\node[rotulo valor, anchor=south, yshift=2pt] at (axis cs:2024,-1.7374) {$-$1,74};
\end{axis}
\end{tikzpicture}

\vspace{8pt}
\begin{tikzpicture}
\begin{axis}[
  estilo ipt,
  width=0.94\linewidth, height=4.6cm,
  title={(b) Grupo de comparação ($T=0$) entre os elegíveis (\%)},
  xmin=2013.6, xmax=2024.4,
  xtick={2014,...,2024},
  xticklabel style={/pgf/number format/1000 sep={}},
  ymin=0, ymax=17,
  ytick={0,4,...,16},
  ymajorgrids,
  xlabel={Ano de nascimento},
  ylabel={$T=0$ (\% dos elegíveis)},
]
\addplot[serie secundaria, mark=diamond*] coordinates {
  (2014,14.7575) (2015,14.6845) (2016,13.3536) (2017,11.8177) (2018,11.3605)
  (2019,11.6874) (2020,11.0875) (2021,9.7911) (2022,9.8718) (2023,9.3829)
  (2024,8.8217)};
\node[rotulo valor, anchor=south, yshift=2pt] at (axis cs:2014,14.7575) {14,76};
\node[rotulo valor, anchor=south, yshift=2pt] at (axis cs:2024,8.8217) {8,82};
\end{axis}
\end{tikzpicture}

\nota{Painel (a): estimativas pontuais de onze modelos anuais independentes, sem intervalo
de confiança anual; diferenças entre anos são descritivas, sem teste formal. Linha
tracejada horizontal: diferença nula. Cada ano tem sua própria população de sobreposição.
Painel (b): denominador igual aos nascimentos elegíveis do ano.}
\fonte{\arquivo{notebooks/principal/07_evolucao_ato_sp_capital_2014_2024.ipynb} (seções 5 e 6).}
\end{figure}

O grupo de comparação diminuiu no período (Figura~\ref{fig:serie}, painel b): de 24.348
nascimentos (14,76\% dos elegíveis) em 2014 para 10.118 (8,82\%) em 2024, queda de
58,44\%, enquanto o grupo $T=1$ caiu 25,64\% (de 140.639 para 104.576) e o total de
elegíveis, 30,48\%. A proporção de $T=0$ teve pequenas reversões em 2019 e 2022 e,
portanto, não caiu de forma estritamente monotônica. Em 2024 estão os menores tamanhos
efetivos dos dois grupos, ponto a considerar em qualquer quantificação futura da precisão.

\subsection{Diagnósticos anuais}

Os diagnósticos repetem o padrão de 2024 (Tabela~\ref{tab:seriediag}). Antes da
ponderação, o maior $|\mathrm{SMD}|$ entre as covariáveis ajustadas variou de 0,452 a
0,548; depois, não passou de 0,0070 (2020), e nenhum termo ajustado ficou acima de 0,1 em
nenhum ano. Como discutido na Seção~\ref{sec:propriedades}, esse balanço é esperado por
construção e não demonstra ausência de confundimento. A situação conjugal, fora do modelo,
manteve maior $|\mathrm{SMD}|$ após a ponderação entre 0,184 e 0,207 --- acima de 0,1 em
todos os anos e acima de 0,2 em 2014 e 2019. A retenção do tamanho efetivo foi de pelo
menos 77,4\% em $T=1$ e 99,2\% em $T=0$; a maior participação individual na soma dos pesos
de um grupo foi de 0,0108\%, e a maior massa de pesos fora do suporte comum observado, de
0,0329\% (2023, $T=1$). Não houve recorte de suporte nem exclusão de registros.

\begin{table}[htbp]
\centering
\caption{Diagnósticos de balanço, escore e pesos por ano, São Paulo (capital), 2014--2024}
\label{tab:seriediag}
\footnotesize
\setlength{\tabcolsep}{3.6pt}
\begin{tabular}{@{}l S[table-format=1.4] S[table-format=1.4] S[table-format=1.4]
  S[table-format=1.4] c S[table-format=2.1] S[table-format=1.4] S[table-format=1.4]@{}}
\toprule
 & \multicolumn{3}{c}{\cab{Maior $|\mathrm{SMD}|$}} & \multicolumn{2}{c}{\cab{Escore}} &
  {\cab{$n_{\mathrm{ef}}/n$}} & \multicolumn{2}{c}{\cab{Peso fora (\%)}} \\
\cmidrule(lr){2-4}\cmidrule(lr){5-6}\cmidrule(lr){7-7}\cmidrule(l){8-9}
\cab{Ano} & {\cab{Antes}} & {\cab{ATO}} & {\cab{Conjugal\textsuperscript{a}}} &
  {\cab{Média}} & \cab{Suporte comum} & {\cab{$T=1$ (\%)}} & {\cab{$T=1$}} & {\cab{$T=0$}} \\
\midrule
2014 & 0.5482 & 0.0022 & 0.2026 & 0.8524 & 0,3240--0,9795 & 77.4 & 0.0044 & 0.0000 \\
2015 & 0.5324 & 0.0010 & 0.1839 & 0.8532 & 0,3428--0,9880 & 77.9 & 0.0000 & 0.0014 \\
2016 & 0.5022 & 0.0010 & 0.1908 & 0.8665 & 0,3938--0,9694 & 79.4 & 0.0102 & 0.0000 \\
2017 & 0.5140 & 0.0027 & 0.1898 & 0.8818 & 0,5020--0,9739 & 79.2 & 0.0082 & 0.0000 \\
2018 & 0.5181 & 0.0036 & 0.1991 & 0.8864 & 0,4415--0,9753 & 78.0 & 0.0006 & 0.0020 \\
2019 & 0.4879 & 0.0030 & 0.2066 & 0.8831 & 0,5907--0,9712 & 79.9 & 0.0132 & 0.0000 \\
2020 & 0.4900 & 0.0070 & 0.1953 & 0.8891 & 0,5014--0,9730 & 80.2 & 0.0039 & 0.0000 \\
2021 & 0.5093 & 0.0016 & 0.1943 & 0.9021 & 0,5557--0,9810 & 78.7 & 0.0317 & 0.0000 \\
2022 & 0.5063 & 0.0044 & 0.1997 & 0.9013 & 0,5875--0,9812 & 78.7 & 0.0000 & 0.0053 \\
2023 & 0.4662 & 0.0043 & 0.1967 & 0.9062 & 0,6555--0,9769 & 80.4 & 0.0329 & 0.0000 \\
\destaque
2024 & 0.4518 & 0.0031 & 0.1945 & 0.9118 & 0,6910--0,9760 & 81.3 & 0.0223 & 0.0000 \\
\bottomrule
\end{tabular}
\nota{Antes e ATO: covariáveis e termos incluídos no escore, sem e com ponderação.
\textsuperscript{a} Situação conjugal, não incluída no modelo, após a ponderação. Média:
média do escore estimado, igual à proporção de $T=1$ no ano. Suporte comum: interseção dos
intervalos observados do escore nos dois grupos. Peso fora: participação dos registros
fora do suporte comum na soma dos pesos do grupo. $n_{\mathrm{ef}}/n$ em $T=0$: pelo menos
99,2\% em todos os anos.}
\fonte{\arquivo{notebooks/principal/07_evolucao_ato_sp_capital_2014_2024.ipynb} (seções 3 e 4).}
\end{table}

\subsection{Leitura da série}

Quatro cautelas orientam a leitura das estimativas anuais.
\begin{itemize}
\item \emph{Cada ano tem sua própria população de sobreposição.} A população-alvo do ATO
é definida pelos pesos e, portanto, pela distribuição das covariáveis e do escore em cada
ano. A média do escore passou de 0,8524 em 2014 para 0,9118 em 2024, e o limite inferior
do suporte comum, de 0,3240 para 0,6910, acompanhando a redução do grupo de comparação.
\item \emph{Os estimandos anuais não são necessariamente comparáveis.} Como a
população-alvo muda, as estimativas não descrevem um efeito sobre uma população fixa;
diferenças entre anos podem refletir mudança de composição, e não apenas da associação.
\item \emph{Não foram estimados intervalos de confiança anuais.} Também não foram testados
contrastes entre anos; nenhuma mudança temporal é apresentada como estatisticamente
significativa. Tamanho efetivo elevado não equivale a precisão demonstrada.
\item \emph{As oscilações não demonstram mudanças causais.} Cada estimativa anual é uma
associação ajustada pelas covariáveis observadas na população de sobreposição do ano. As
oscilações históricas não demonstram mudanças causais no efeito do início do pré-natal e
não podem ser atribuídas a ele nem a eventos externos.
\end{itemize}
A série amplia a descrição da consistência da direção da associação ao longo do tempo,
mas não elimina confundimento não observado, seleção por nascido vivo nem os problemas de
temporalidade da exposição e das covariáveis (Seções~\ref{sec:selecao}
e~\ref{sec:discussao}).

% =====================================================================
\section{Temporalidade do início do pré-natal e seleção da população}\label{sec:selecao}

Esta seção descreve, para 2024, a relação entre a duração da gestação, a oportunidade de
iniciar o pré-natal em meses mais tardios e a composição da população analítica. A análise
é exclusivamente descritiva: não houve nova estimação, o estimando não muda e a duração
gestacional não entra no ajuste.

\subsection{Oportunidade temporal de iniciar o pré-natal}

Uma gestação que termina cedo não chega aos meses mais tardios e, por isso, não pode ter
registrado nesses meses o início do pré-natal. A possibilidade de pertencer ao grupo de
comparação depende, assim, em parte da própria duração da gestação, que por sua vez está
fortemente ligada ao baixo peso. A restrição existe por construção. Em uma formulação
anterior do projeto, com comparador nos meses 4--9, ela alcançava todas as gestações
pré-termo; a restrição aos meses 4--6 reduz a janela de início tardio considerada, mas não
demonstra que o viés diminua nem indica sua direção.

O comparador 4--6 continua sujeito ao problema, sobretudo pelo mês 6, nas gestações que
terminam antes do fim do sexto mês. Nas faixas oficiais de duração do campo
\campo{GESTACAO}, isso abrange as gestações com menos de 22 semanas e parte das de 22 a 27
semanas, faixa em que cai, pelas convenções usuais, a fronteira do sexto mês --- sem que
se fixe a semana exata. Juntas, essas duas faixas somam 546 nascimentos da população
analítica (0,48\%), 59 deles no grupo $T=0$. A partir de 28 semanas, a oportunidade de
iniciar o pré-natal em qualquer mês de 1 a 6 está completa.

\subsection{Duração da gestação na população analítica}

A duração gestacional está informada em quase toda a população analítica: apenas 7
registros não têm \campo{SEMAGESTAC} nem \campo{GESTACAO}, e os dois campos concordam em
todos os registros preenchidos. Essa concordância indica consistência entre campos, e não
acurácia, pois ambos podem derivar da mesma estimativa. A validade da medida não foi
avaliada: em 2024, 37,62\% das idades gestacionais foram estimadas por exame físico,
método menos preciso, e 18,68\% trazem no campo de método de estimativa
(\campo{TPMETESTIM}) o código 8, ausente do dicionário e não interpretado. Erros de uma ou
duas semanas perto das fronteiras das faixas são plausíveis; as faixas servem para
descrição aproximada, não como medida validada.

A distribuição concentra-se a termo: 89,96\% dos nascimentos tiveram de 37 a 41 semanas
(Tabela~\ref{tab:duracao}). O baixo peso é quase universal abaixo de 32 semanas, atinge
48,02\% de 32 a 36 semanas e fica em 3,72\% de 37 a 41. Os 546 nascimentos com menos de 28
semanas (0,48\% da população) concentram cerca de 5\% dos casos de baixo peso: qualquer
desequilíbrio do comparador nessas faixas pesa no desfecho de forma desproporcional ao seu
tamanho.

\begin{table}[htbp]
\centering
\caption{Duração da gestação, início do pré-natal nos meses 4--6 e baixo peso ao nascer na
população analítica, São Paulo (capital), 2024}
\label{tab:duracao}
\small
\begin{tabular}{@{}l S[table-format=6.0] S[table-format=3.2] S[table-format=2.2]
  S[table-format=4.0] S[table-format=2.2] S[table-format=3.2]@{}}
\toprule
 & \multicolumn{2}{c}{\cab{Nascimentos}} & {\cab{Início 4--6}} &
  \multicolumn{3}{c}{\cab{Baixo peso ao nascer}} \\
\cmidrule(lr){2-3}\cmidrule(lr){4-4}\cmidrule(l){5-7}
\cab{Duração (semanas)} & {\cab{$n$}} & {\cab{\% do total}} & {\cab{\% da faixa}} &
  {\cab{Casos}} & {\cab{\% da faixa}} & {\cab{\% dos casos}} \\
\midrule
$<$ 22     & 15     & 0.01  & 6.67  & 9    & 60.00 & 0.09  \\
22--27     & 531    & 0.46  & 10.92 & 504  & 94.92 & 5.12  \\
28--31     & 981    & 0.86  & 11.52 & 904  & 92.15 & 9.18  \\
32--36     & 9525   & 8.30  & 9.96  & 4574 & 48.02 & 46.44 \\
37--41     & 103180 & 89.96 & 8.65  & 3843 & 3.72  & 39.02 \\
42 ou mais & 455    & 0.40  & 14.29 & 16   & 3.52  & 0.16  \\
Ignorada   & 7      & 0.01  & 28.57 & 0    & 0.00  & 0.00  \\
\midrule
Total      & 114694 & 100.00 & 8.82 & 9850 & 8.59  & 100.00 \\
\bottomrule
\end{tabular}
\nota{Faixas oficiais do campo \campo{GESTACAO}, registradas no parto. Início 4--6:
proporção dos nascimentos da faixa com início do pré-natal nos meses 4 a 6 (grupo $T=0$).
Os percentuais da linha ``Ignorada'' referem-se a 7 registros. A tabela descreve a
população já selecionada por início nos meses 1--6; diferenças dentro das faixas não são
efeitos causais.}
\fonte{\arquivo{notebooks/principal/08_temporalidade_prenatal_sp_capital_2024.ipynb} (seções 2 a 4).}
\end{table}

Dentro da população analítica, a composição do comparador não mostra o padrão que o
truncamento produziria sozinho (Figura~\ref{fig:duracao}, painel a). O início nos meses
4--6 está presente em todas as faixas a partir de 22 semanas, com proporções iguais ou
maiores que a termo: 10,92\% dos nascimentos de 22 a 27 semanas, contra 8,65\% dos de 37 a
41. A única faixa sem inícios nos meses 5 e 6 é a de menos de 22 semanas, com 15
nascimentos, número pequeno demais para ser informativo. Por grupo, a proporção de
nascimentos com menos de 28 semanas foi de 0,47\% em $T=1$ e 0,58\% em $T=0$ (0,80\% no
mês 6), e a de nascimentos de 28 a 36 semanas, de 9,03\% e 10,50\%. O risco bruto de baixo
peso cresce do mês 1 (7,9\%) ao mês 6 (11,4\%), acompanhado pela fração de nascimentos
pré-termo, de 9,3\% no mês 1 a 12,5\% e 11,6\% nos meses 5 e 6 (Figura~\ref{fig:duracao},
painel b). A duração gestacional, portanto, difere entre os grupos que formam o contraste.

\begin{figure}[htbp]
\centering
\caption{Início do pré-natal nos meses 4--6 segundo a duração da gestação e composição do
risco bruto de baixo peso por mês de início, São Paulo (capital), 2024}
\label{fig:duracao}
% Figura 5 (Etapa 2) — Temporalidade do início do pré-natal, São Paulo (capital), 2024.
% Painel (a): % dos nascimentos de cada faixa oficial de GESTACAO com início nos meses 4, 5
%   e 6, copiado do gráfico Plotly salvo em
%   notebooks/principal/08_temporalidade_prenatal_sp_capital_2024.ipynb (seção 3);
%   exclui os 7 registros sem duração.
% Painel (b): risco bruto de baixo peso por mês de início, separado em nascimentos
%   pré-termo (< 37 semanas) e a termo, do gráfico Plotly do mesmo notebook (seção 4);
%   denominador: nascimentos do mês com duração conhecida. Decomposição contábil, sem ajuste.
% Os dois painéis são alinhados pela base da área de plotagem (baseline=(eixo.south)).
\begin{minipage}[b]{0.49\linewidth}
\centering
\begin{tikzpicture}[baseline=(eixo.south)]
\begin{axis}[
  name=eixo,
  estilo ipt,
  scale only axis,
  width=5.0cm, height=4.8cm,
  title={(a) Início nos meses 4–6, por duração},
  ybar stacked,
  bar width=11pt,
  xmin=0.4, xmax=6.6,
  xtick={1,...,6},
  xticklabels={$<$22, 22–27, 28–31, 32–36, 37–41, $\geq$42},
  xticklabel style={font=\sffamily\scriptsize},
  ymin=0, ymax=19,
  ytick={0,5,10,15},
  ymajorgrids,
  xlabel={Duração da gestação (semanas)},
  ylabel={\% dos nascimentos da faixa},
  legend style={at={(0.5,0.98)}, anchor=north, legend columns=3,
    /tikz/every even column/.append style={column sep=4pt}},
]
\addplot[fill=iptazul!30, draw=white, line width=0.3pt] coordinates {
  (1,6.6667) (2,4.5198) (3,5.3007) (4,4.9029) (5,4.6395) (6,7.4725)};
\addlegendentry{Mês 4}
\addplot[fill=iptazul!60, draw=white, line width=0.3pt] coordinates {
  (1,0) (2,3.7665) (3,3.9755) (4,3.3176) (5,2.5267) (6,4.8352)};
\addlegendentry{Mês 5}
\addplot[fill=iptazul, draw=white, line width=0.3pt] coordinates {
  (1,0) (2,2.6365) (3,2.2426) (4,1.7428) (5,1.4887) (6,1.9780)};
\addlegendentry{Mês 6}
% Total dos meses 4–6 em cada faixa
\node[rotulo valor, anchor=south] at (axis cs:1,6.6667) {6,7};
\node[rotulo valor, anchor=south] at (axis cs:2,10.9228) {10,9};
\node[rotulo valor, anchor=south] at (axis cs:3,11.5188) {11,5};
\node[rotulo valor, anchor=south] at (axis cs:4,9.9633) {10,0};
\node[rotulo valor, anchor=south] at (axis cs:5,8.6549) {8,7};
\node[rotulo valor, anchor=south] at (axis cs:6,14.2857) {14,3};
\end{axis}
\end{tikzpicture}
\end{minipage}%
\hfill
\begin{minipage}[b]{0.49\linewidth}
\centering
\begin{tikzpicture}[baseline=(eixo.south)]
\begin{axis}[
  name=eixo,
  estilo ipt,
  scale only axis,
  width=5.0cm, height=4.8cm,
  title={(b) Baixo peso por mês de início},
  ybar stacked,
  bar width=11pt,
  xmin=0.4, xmax=6.6,
  xtick={1,...,6},
  ymin=0, ymax=16,
  ytick={0,4,...,16},
  ymajorgrids,
  xlabel={Mês de início do pré-natal},
  ylabel={\% dos nascimentos do mês},
  legend style={at={(0.02,0.98)}, anchor=north west},
]
% Linha divisória: meses 4–6 à direita (grupo T=0)
\draw[densely dashed, iptcinza!75, line width=0.7pt] (axis cs:3.5,0) -- (axis cs:3.5,16);
\addplot[fill=iptazul, draw=white, line width=0.3pt] coordinates {
  (1,4.8887) (2,5.1120) (3,5.6618) (4,5.5918) (5,7.1238) (6,7.3268)};
\addlegendentry{Pré-termo ($<$ 37 sem.)}
\addplot[fill=iptazul!35, draw=white, line width=0.3pt] coordinates {
  (1,3.0462) (2,3.4129) (3,3.4682) (4,3.9143) (5,4.2943) (6,4.1214)};
\addlegendentry{A termo ($\geq$ 37 sem.)}
% Risco bruto total do mês
\node[rotulo valor, anchor=south] at (axis cs:1,7.9349) {7,9};
\node[rotulo valor, anchor=south] at (axis cs:2,8.5249) {8,5};
\node[rotulo valor, anchor=south] at (axis cs:3,9.1300) {9,1};
\node[rotulo valor, anchor=south] at (axis cs:4,9.5061) {9,5};
\node[rotulo valor, anchor=south] at (axis cs:5,11.4181) {11,4};
\node[rotulo valor, anchor=south] at (axis cs:6,11.4482) {11,4};
\end{axis}
\end{tikzpicture}
\end{minipage}

\nota{Painel (a): percentual dos nascimentos de cada faixa oficial de duração
(\campo{GESTACAO}) com início nos meses 4, 5 e 6; $n$ por faixa: 15 ($<$ 22 semanas), 531,
981, 9.525, 103.180 e 455 ($\geq$ 42). Painel (b): percentual dos nascimentos de cada mês
de início com duração conhecida ($n$ de 35.224 no mês 1 a 1.747 no mês 6), separado em
nascimentos pré-termo e a termo; linha divisória tracejada: à direita, meses do grupo $T=0$. Decomposição
contábil do risco bruto, sem ajuste; não constitui análise de mediação.}
\fonte{\arquivo{notebooks/principal/08_temporalidade_prenatal_sp_capital_2024.ipynb} (seções 3 e 4).}
\end{figure}

Essa composição não afasta o problema. Ela mistura o truncamento, diferenças de perfil de
risco entre quem inicia o pré-natal cedo e quem o inicia tarde e um possível efeito do
próprio pré-natal sobre a duração da gestação, e refere-se apenas aos registros que já
passaram pelo filtro de início nos meses 1--6. A proporção de 0,48\% descreve essa
população selecionada: não demonstra que o viés seja pequeno nem o limita a essa fração.

\subsection{Registros fora da população: início ausente, ignorado ou tardio}

Antes do filtro de início do pré-natal, mantidos os demais critérios da população
analítica --- residência materna na capital, nascimento em 2024, gestação única e peso de
500 a 6.000\,g ---, havia 117.422 nascimentos. Desses, 114.694 (97,68\%) foram incluídos e
2.728 (2,32\%) ficaram fora: 1.735 com início nos meses 7 a 9, os mesmos da
Tabela~\ref{tab:fluxo}, e 993 sem informação de início, dos quais 889 com o campo ausente e
104 com o código 99 (ignorado); nenhum outro código foi observado
(Tabela~\ref{tab:selecao}).

\begin{table}[htbp]
\centering
\caption{Nascimentos antes do filtro de início do pré-natal, por informação de início e
duração da gestação, São Paulo (capital), 2024}
\label{tab:selecao}
\small
\begin{tabularx}{\linewidth}{@{}L S[table-format=6.0] S[table-format=3.2]
  S[table-format=3.0] S[table-format=1.2] S[table-format=2.2] S[table-format=2.2]
  S[table-format=2.0]@{}}
\toprule
 & \multicolumn{2}{c}{\cab{Nascimentos}} & \multicolumn{2}{c}{\cab{$<$ 28 semanas}} &
  {\cab{28--36}} & {\cab{37 ou mais}} & {\cab{Duração}} \\
\cmidrule(lr){2-3}\cmidrule(lr){4-5}\cmidrule(lr){6-6}\cmidrule(lr){7-7}\cmidrule(l){8-8}
\cab{Início do pré-natal} & {\cab{$n$}} & {\cab{\% do total}} & {\cab{$n$}} & {\cab{\%}} &
  {\cab{\%}} & {\cab{\%}} & {\cab{desconhecida}} \\
\midrule
Meses 1--3 ($T=1$)     & 104576 & 89.06  & 487 & 0.47 & 9.03  & 90.50 & 5  \\
Meses 4--6 ($T=0$)     & 10118  & 8.62   & 59  & 0.58 & 10.50 & 88.90 & 2  \\
Meses 7--9             & 1735   & 1.48   & 3   & 0.17 & 10.32 & 89.45 & 1  \\
Ausente ou ignorado    & 993    & 0.85   & 40  & 4.03 & 21.85 & 73.92 & 2  \\
\addlinespace
\destaque
Incluídos (meses 1--6) & 114694 & 97.68  & 546 & 0.48 & 9.16  & 90.36 & 7  \\
Excluídos              & 2728   & 2.32   & 43  & 1.58 & 14.52 & 83.80 & 3  \\
\midrule
Total antes do filtro  & 117422 & 100.00 & 589 & 0.50 & 9.28  & 90.21 & 10 \\
\bottomrule
\end{tabularx}
\nota{Recorte: residência materna no município de São Paulo, nascimento em 2024, gestação
única e peso de 500 a 6.000\,g, sem filtro de \campo{MESPRENAT}. Ausente ou ignorado: 889
registros com campo ausente e 104 com código 99. Duração pelas faixas oficiais de
\campo{GESTACAO} (28--36: códigos 3 e 4; 37 ou mais: códigos 5 e 6). Percentuais de
duração sobre o total da linha, incluindo duração desconhecida. As quatro primeiras linhas
são categorias exclusivas; as três últimas são subtotais e total.}
\fonte{\arquivo{notebooks/principal/08_temporalidade_prenatal_sp_capital_2024.ipynb} (seção 5.1).}
\end{table}

Há diferenças de composição entre incluídos e excluídos. A proporção de gestações com
menos de 28 semanas foi de 1,58\% entre os excluídos (43 de 2.728), contra 0,48\% entre os
incluídos (546 de 114.694). Essa concentração está nos registros sem informação de início
(4,03\%, 40 de 993), e não nos inícios nos meses 7--9 (0,17\%, 3 de 1.735). A duração é
desconhecida em 3 excluídos e 7 incluídos, mantidos nos denominadores.

Informação de início ausente ou ignorada não equivale a ausência de pré-natal. Um
mecanismo possível, que os dados não permitem verificar, é o de uma gestação que terminou
antes do mês em que o pré-natal começaria ser registrada sem pré-natal, com início ignorado
ou com início fora dos meses 1--6: nesse caso, ela sai da população por seleção, e não muda
de grupo. A maior concentração de gestações curtas entre os excluídos descreve seleção
observada anterior à comparação; não identifica a magnitude nem a direção do possível viés
causal e não demonstra, por si só, que ele seja grande.

\subsection{Restrição aos nascidos vivos}

A população de partida também é selecionada pela condição de nascido vivo. O SINASC
registra apenas nascidos vivos, e este estudo não vinculou seus registros aos do Sistema
de Informações sobre Mortalidade (SIM). Gestações que terminaram em óbito fetal não são
observadas nem na população analítica nem na população ampliada desta seção, de modo que
a seleção por nascido vivo não pode ser avaliada com os dados atuais
(Seção~\ref{sec:selecao-nascido-vivo}).

\subsection{O que a análise permite e o que não permite concluir}

Os dados permitem descrever a disponibilidade da duração gestacional, a composição de cada
mês de início, a concentração do baixo peso nas durações curtas e as diferenças de
composição entre registros incluídos e excluídos. Não permitem estimar o tamanho nem o
sinal de um viés de truncamento ou de seleção, nem afirmar que ele seja pequeno, nem ler
diferenças dentro de cada faixa de duração como efeitos causais. Também não justificam
acrescentar a duração gestacional ao ajuste: registrada no parto, ela pode ser mediadora
ou variável de seleção \cite{wilcox2011,schisterman2009}. A duração gestacional permanece
fora do ajuste principal, e a estimativa da Seção~\ref{sec:resultados} não foi alterada.

Uma comparação restrita aos meses 4--5 retiraria do grupo $T=0$ os 1.747 nascimentos com
início no mês 6 (17\% do grupo), dos quais 14 com menos de 28 semanas; a composição do mês
6 por duração é próxima à dos meses 4--5. Esses números descrevem o que tal análise
removeria, mas não permitem afirmar que ela corrigiria o problema: continuaria restrita aos
registros incluídos, não alcançaria a seleção anterior ao filtro de início e misturaria
perda de amostra, perfil de quem inicia no mês 6 e eventual efeito do truncamento. Se
realizada, deverá ser lida como análise de sensibilidade, e não como correção. Ela não foi
executada.

% =====================================================================
\section{Discussão e limitações}\label{sec:discussao}

\subsection{O que foi observado e o que a ponderação ajustou}

Entre os nascimentos de mães residentes na capital paulista em 2024, o risco de baixo peso
foi menor no grupo com início do pré-natal nos meses 1--3. A diferença bruta de
$-2{,}00$\pp{} reduziu-se a $-1{,}74$\pp{} depois que as distribuições de idade,
escolaridade, raça/cor, paridade e perdas fetais anteriores foram equilibradas entre os
grupos, e o intervalo de confiança da diferença ajustada não incluiu zero. Na análise de
sensibilidade, a inclusão da situação conjugal levou a $-1{,}63$\pp, na mesma direção.
Repetida em cada ano de 2014 a 2024, a mesma especificação produziu diferenças ajustadas
negativas nos onze anos, entre $-0{,}58$ e $-1{,}90$\pp, sem tendência monotônica, sem
intervalos de confiança anuais e em populações de sobreposição que mudam de um ano para
outro (Seção~\ref{sec:serie}).

O que a ponderação fez foi tornar comparáveis, em média, as distribuições das
características observadas incluídas no modelo, definindo uma população ponderada de
sobreposição. Ela não pareou gestantes individualmente, não equilibrou características não
registradas e não equilibrou completamente a situação conjugal, deixada fora do modelo
principal. A redução de cerca de 13\% na magnitude da diferença indica que parte da
diferença bruta acompanhava a composição dos grupos nas covariáveis observadas; a parte
remanescente não pode ser atribuída ao momento de início sem as hipóteses discutidas a
seguir.

\subsection{O que permanece sem solução}

A principal ameaça é o confundimento não observado. Renda, cobertura por plano privado ou
uso exclusivo do SUS, planejamento e desejo da gestação, tabagismo, estado nutricional,
condições clínicas prévias e acesso geográfico à unidade de saúde não constam do SINASC e
podem afetar tanto o momento de início do pré-natal quanto o peso ao nascer
(Figura~\ref{fig:dag}). A escolaridade capta parte do gradiente socioeconômico, mas não o
esgota. A AUC modesta do escore indica que boa parte da variação no momento de início não
é explicada pelas covariáveis registradas; o balanço obtido nada informa sobre essa parte.

\begin{figure}[htbp]
\centering
\caption{Diagrama causal conceitual adotado como hipótese de trabalho}
\label{fig:dag}
% Diagrama causal conceitual (DAG de trabalho); Figura 6 na Etapa 2 (Figura 4 na Etapa 1).
% Elaborado a partir de docs/methodology/DAG_INICIAL.md e das limitações registradas
% em docs/sp_capital_ato/RESUMO_PARA_REDACAO.md. As setas representam hipóteses de trabalho,
% não relações estimadas. Contorno contínuo: observado; tracejado: não observado;
% contorno duplo: seleção pela condição de nascido vivo.
\begin{tikzpicture}[
  font=\sffamily\footnotesize,
  >={Stealth[length=2.4mm, width=1.8mm]},
  no/.style={draw=iptazul, line width=0.8pt, rounded corners=2pt, align=center,
    inner sep=5pt, fill=white, text=iptcinza!25!black,
    execute at begin node={\hyphenpenalty=10000\exhyphenpenalty=10000}},
  nao observado/.style={no, draw=iptcinza!80, dashed},
  posterior/.style={no, draw=iptcinza!80},
  selecao/.style={no, double=white, double distance=1.3pt},
  seta/.style={->, line width=0.85pt, iptazul},
  seta fraca/.style={->, line width=0.85pt, iptcinza!80},
  seta hipotetica/.style={->, line width=0.85pt, iptcinza!80, dashed},
  rotulo/.style={font=\sffamily\scriptsize, text=iptcinza!40!black, fill=white, inner sep=1.5pt},
]
\node[no, text width=3.4cm] (T) at (0,0)
  {\textbf{T} --- início do pré-natal\\meses 1–3 \emph{versus} 4–6};
\node[no, text width=3.9cm] (Y) at (10.6,0)
  {\textbf{Y} --- baixo peso ao nascer\\(peso $<$ 2.500 g)};
\node[no, text width=9.4cm] (X) at (5.3,2.55)
  {\textbf{X} --- características maternas observadas e ajustadas\\
   idade, escolaridade, raça/cor, paridade, perdas fetais anteriores};
\node[posterior, text width=5.2cm] (M) at (5.3,-1.05)
  {\textbf{M} --- assistência e curso da gestação\\
   consultas e idade gestacional\\(posteriores a T; não ajustados)};
\node[nao observado, text width=6.6cm] (U) at (7.3,-3.55)
  {\textbf{U} --- fatores não registrados no SINASC\\
   renda, plano privado ou SUS, planejamento da gestação, tabagismo,
   estado nutricional, condições clínicas prévias, acesso à unidade de saúde};
\node[selecao, text width=3.7cm] (C) at (0,-3.55)
  {\textbf{C} --- nascimento vivo\\análise restrita a $C\!=\!1$};

\draw[seta] ($(X.south west)+(0.9,0)$) -- (T.north);
\draw[seta] ($(X.south east)+(-0.9,0)$) -- (Y.north);
\draw[seta, line width=1.6pt] (T.east) -- node[rotulo, above=2pt] {relação investigada} (Y.west);
\draw[seta fraca] (T.south east) -- (M.west);
\draw[seta fraca] (M.east) -- (Y.south west);
\draw[seta hipotetica] ($(U.north west)+(0.4,0)$) -- ($(T.south)+(0.6,0)$);
\draw[seta hipotetica] ($(U.north east)+(-0.7,0)$) -- (Y.south);
\draw[seta hipotetica] ($(T.south)+(-0.5,0)$) -- node[rotulo, left=2pt, pos=0.5] {possível}
  ($(C.north)+(-0.5,0)$);
\end{tikzpicture}

\nota{Contorno contínuo azul: variáveis observadas; contorno duplo: seleção pela condição
de nascido vivo; contorno cinza: variáveis posteriores à exposição, não ajustadas;
contorno e setas tracejados: fatores não observados e mecanismos possíveis, não avaliados
com os dados. Todas as setas representam hipóteses de trabalho, não relações estimadas.}
\fonte{elaboração própria a partir de \arquivo{docs/methodology/DAG_INICIAL.md} e das
limitações registradas em \arquivo{docs/sp_capital_ato/RESUMO_PARA_REDACAO.md}.}
\end{figure}

A situação conjugal ilustra o dilema entre confundimento e temporalidade. Ela foi excluída
do modelo principal porque pode mudar depois do início do pré-natal; incluída na
sensibilidade, alterou pouco a estimativa, mas também mudou a população-alvo. A decisão
sobre seu papel deve ser tomada com base no diagrama causal, e não na magnitude da
mudança da estimativa.

\subsection{Temporalidade da exposição e duração gestacional}\label{sec:temporalidade}

O mês de início do pré-natal é registrado na Declaração de Nascido Vivo, no parto, a
partir de recordação ou de transcrição do cartão da gestante. Erros de classificação entre
meses adjacentes, como 3 e 4, são plausíveis e não foram quantificados.

Há também uma dificuldade estrutural. Gestações mais curtas têm menos tempo para que o
pré-natal comece tardiamente, e a prematuridade está fortemente ligada ao baixo peso.
Restringir o grupo de comparação aos meses 4--6 reduz a janela de início tardio
considerada, mas não foi demonstrado quanto isso altera a ligação entre duração da
gestação e classificação da exposição, nem o viés resultante ou sua direção. A análise
descritiva da Seção~\ref{sec:selecao} mostrou que a duração gestacional difere entre os
grupos e que gestações curtas são mais frequentes entre os registros excluídos por falta
de informação de início; esses padrões descrevem composição e seleção observada, mas não
identificam a magnitude nem a direção do possível viés. Ajustar pela idade gestacional ao
nascer não resolve o problema: trata-se de variável posterior à exposição, e o ajuste por
ela pode introduzir viés \cite{wilcox2011,schisterman2009}. As covariáveis, por sua vez, também
são registradas no parto: a escolaridade pode mudar durante a gestação em adolescentes, e
o campo de perdas fetais anteriores mistura perdas espontâneas e abortos, com provável
sub-registro.

\subsection{Seleção por nascimento vivo}\label{sec:selecao-nascido-vivo}

A população inclui apenas gestações que resultaram em nascido vivo. Se o momento de início
do pré-natal influenciar perdas fetais, e se houver fatores que afetem tanto a
sobrevivência fetal quanto o peso ao nascer, a comparação restrita a nascidos vivos pode
ser enviesada \cite{snowden2018}. O SINASC registra apenas nascidos vivos, e este estudo
não vinculou seus registros aos do Sistema de Informações sobre Mortalidade (SIM).
Desfechos gestacionais fora desse universo, como os óbitos fetais, não são observados na
população analítica, o que impede avaliar essa seleção com os dados atuais. Essa seleção
soma-se à observada antes do filtro de início do pré-natal (Seção~\ref{sec:selecao}), que
se refere apenas a nascidos vivos registrados.

\subsection{Alcance da população-alvo e da inferência}

O estimando ATO refere-se a uma população definida pelos próprios pesos --- mais jovem e
menos escolarizada que o total --- e não se generaliza automaticamente para todas as
gestantes da capital (ATE), para as que iniciaram cedo (ATT), para outros municípios ou
para outros anos. Mudar as covariáveis do escore muda essa população, como mostrou a
análise de sensibilidade. A série anual da Seção~\ref{sec:serie} não supre essa lacuna:
cada ano tem sua própria população de sobreposição, e as estimativas anuais não descrevem
um efeito sobre uma população fixa. Pela mesma razão, as estimativas da fase anterior do
projeto, de âmbito nacional e obtidas com AIPW e erros-padrão agrupados por município
(por exemplo, $-1{,}48$\pp{} em 2024), referem-se a outra população, outro estimando e
outro estimador e não devem ser combinadas com as deste relatório.

A validade da ponderação depende ainda da especificação do escore: um modelo logístico de
efeitos principais com idade quadrática pode não representar interações relevantes entre
as covariáveis, e o estimador não é duplamente robusto. Quanto à inferência, o bootstrap
tratou os nascimentos como independentes; o intervalo apresentado pode ser estreito demais
na presença de agrupamento e não cobre vieses sistemáticos.

\subsection{O que não se pode concluir}

Os resultados não permitem afirmar que iniciar o pré-natal nos meses 1--3 reduz o risco de
baixo peso em 1,74\pp, nem que antecipar o início para gestantes que hoje o iniciam nos
meses 4--6 produziria essa redução. Também não autorizam extrapolar a magnitude para a
população total de gestantes, para outros anos ou para o estado, nem afirmar que a
associação tenha mudado entre 2014 e 2024: as oscilações anuais não foram testadas e não
podem ser atribuídas ao início do pré-natal nem a eventos externos. A análise de
temporalidade não permite afirmar que o viés ligado à duração gestacional seja pequeno,
grande ou de determinada direção. O que os dados sustentam
é uma diferença de risco persistente após o equilíbrio das características maternas
observadas, em uma população-alvo bem definida, cuja interpretação causal depende de
hipóteses não verificáveis com o SINASC.

\subsection{Questões em aberto}

As questões a seguir não foram abordadas nesta etapa.
\begin{enumerate}
\item Robustez da associação a confundimento não medido, por exemplo com o valor E
\cite{vanderweele2017}, e busca de um controle negativo defensável.
\item Estimação da incerteza com dependência entre observações, por bootstrap ou
variância agrupada por estabelecimento de nascimento ou região de residência.
\item Intervalo de confiança para a sensibilidade com situação conjugal, necessário caso
ela seja usada além de forma descritiva.
\item Incerteza das estimativas anuais e dos contrastes entre anos, considerando a redução
do grupo $T=0$ e a mudança da população de sobreposição, condição para afirmar qualquer
mudança temporal.
\item Revisão do conjunto de ajuste com base no diagrama causal, incluindo a situação
conjugal e campos possivelmente anteriores à exposição ainda não usados.
\item Desenhos que reduzam a dependência entre duração gestacional e oportunidade de
início tardio, sem condicionar em variáveis posteriores à exposição; uma comparação
restrita aos meses 4--5 seria uma análise de sensibilidade, não uma correção
(Seção~\ref{sec:selecao}).
\item Vinculação individual entre SINASC e SIM, incluindo óbitos fetais, para examinar a
seleção por nascido vivo, se houver autorização e base adequada.
\item Extensão da especificação ao estado de São Paulo, cuja população de sobreposição
seria diferente.
\end{enumerate}

% =====================================================================
{\small\setstretch{1.0}\bibliography{referencias}}

\end{document}
